\documentclass[UTF8,12pt,a4paper,oneside]{ctexbook}
\usepackage[left=3.0cm,right=2.0cm,top=3.0cm,bottom=2.5cm]{geometry}
\usepackage{setspace}
\usepackage{indentfirst}
\usepackage{hyperref}
\hypersetup{hidelinks}
\usepackage{booktabs}
\usepackage{array}
\usepackage{longtable}
\usepackage{tocloft}
\usepackage{graphicx}
\usepackage{pdfpages}
\usepackage{amsmath}
\usepackage{float}
\usepackage{needspace}
\usepackage{tabularx}
\usepackage{caption}
\usepackage{fancyhdr}
\usepackage{makecell}
\usepackage{titlesec}
\usepackage{chngcntr}
\usepackage{etoolbox}
\usepackage{tikz}
\usetikzlibrary{arrows.meta,positioning,shadows.blur,shapes.geometric,calc}
\usepackage[
  backend=biber,
  style=gb7714-2015,
  maxnames=3,
  minnames=3
]{biblatex}
\addbibresource{refs.bib}
\AtBeginBibliography{\sloppy}
\setstretch{1.5}
\setmainfont{Times New Roman}
\setCJKmainfont{SimSun}[BoldFont=SimSun,BoldFeatures={FakeBold=2}]
\setCJKsansfont{Microsoft YaHei}
\setCJKmonofont{FangSong}
% 使用 ctex 内置黑体命令，避免系统字体兼容问题
\setCJKfamilyfont{song}{SimSun}[BoldFont=SimSun,BoldFeatures={FakeBold=2}]
\xeCJKsetup{AutoFakeBold=true}
\ctexset{
  chapter={
    name={,},
    number=\arabic{chapter},
    format=\centering\heiti\bfseries\zihao{3},
    aftername=\hspace{0.5em},
    beforeskip=24pt,
    afterskip=18pt
  },
  section={
    format=\heiti\bfseries\zihao{-4},
    aftername=\hspace{0.5em},
    beforeskip=12pt,
    afterskip=6pt
  },
  subsection={
    format=\songti\bfseries\zihao{-4},
    aftername=\hspace{0.5em},
    beforeskip=6pt,
    afterskip=3pt
  }
}
\setcounter{secnumdepth}{2}
\setcounter{tocdepth}{2}
\setlength{\cftchapindent}{0em}
\setlength{\cftsecindent}{2em}
\setlength{\cftsubsecindent}{4em}
\setlength{\cftchapnumwidth}{2.2em}
\setlength{\cftsecnumwidth}{3.0em}
\setlength{\cftsubsecnumwidth}{4.0em}
\renewcommand{\cftsecpresnum}{\hspace{1em}}
\renewcommand{\cftsubsecpresnum}{\hspace{2em}}
\renewcommand{\cftdotsep}{1}
\DeclareCaptionFont{songwu}{\songti\zihao{5}}
\captionsetup[table]{font={songwu},labelfont={bf},textfont={bf},labelsep=space,justification=centering,singlelinecheck=true}
\captionsetup[figure]{font={songwu},labelfont={bf},textfont={bf},labelsep=space,justification=centering,singlelinecheck=true}
\captionsetup{compatibility=false}
\setlength{\abovecaptionskip}{0.5\baselineskip}
\setlength{\belowcaptionskip}{0.5\baselineskip}
\setlength{\headheight}{15pt}
\setlength{\emergencystretch}{8em}
\setlength{\parindent}{2em}
\setlength{\parskip}{0pt}
\clubpenalty=10000
\widowpenalty=10000
\displaywidowpenalty=10000
\fancyhf{}
\fancyhead[C]{\songti\zihao{-5} 广东工业大学本科毕业设计（论文）}
\fancyfoot[C]{\songti\zihao{-5}\thepage}
\renewcommand{\headrulewidth}{0pt}
\renewcommand{\footrulewidth}{0pt}
% 章节起始页（plain）也使用相同页眉页脚
\fancypagestyle{plain}{%
  \fancyhf{}%
  \fancyhead[C]{\songti\zihao{-5} 广东工业大学本科毕业设计（论文）}%
  \fancyfoot[C]{\songti\zihao{-5}\thepage}%
  \renewcommand{\headrulewidth}{0pt}%
  \renewcommand{\footrulewidth}{0pt}%
}
\renewcommand{\arraystretch}{1.15}
\raggedbottom
\AtBeginDocument{\zihao{-4}}

% 标题防孤行：避免标题落页底、内容到下一页
\pretocmd{\chapter}{\Needspace{6\baselineskip}}{}{}
\pretocmd{\section}{\Needspace{5\baselineskip}}{}{}
\pretocmd{\subsection}{\Needspace{4\baselineskip}}{}{}
% 正文层级“款、项”格式：1、 和 (1)
\newcommand{\kuan}[1]{\par\noindent\hspace*{2em}\CJKfamily{song}\zihao{-4}#1、}
\newcommand{\xiang}[1]{\par\noindent\hspace*{2em}\CJKfamily{song}\zihao{-4}(#1)}
\newcommand{\yanjiuitem}[2]{\par\noindent{\songti\zihao{-4}（{\rmfamily\zihao{-4}#1}）#2}}
\numberwithin{equation}{chapter}
\counterwithin{figure}{chapter}
\counterwithin{table}{chapter}
\renewcommand{\thefigure}{\thechapter.\arabic{figure}}
\renewcommand{\thetable}{\thechapter.\arabic{table}}
\newcommand{\figref}[1]{图~\ref{#1}}
\newcommand{\tabref}[1]{表~\ref{#1}}
\newcommand{\eqreftext}[1]{见式~\eqref{#1}}
\newcommand{\captab}[1]{表~#1}
\newcommand{\capfig}[1]{图~#1}
\tikzset{
  flowblock/.style={
    rounded corners=3pt,
    draw=blue!60!black,
    fill=blue!6,
    line width=0.75pt,
    minimum width=29mm,
    minimum height=7.8mm,
    align=center,
    font=\zihao{-5},
    blur shadow={shadow blur steps=4, shadow blur extra rounding=1.2pt, shadow xshift=0.8pt, shadow yshift=-0.8pt, shadow opacity=16}
  },
  flowdecision/.style={
    diamond,
    draw=blue!60!black,
    fill=cyan!8,
    line width=0.75pt,
    aspect=1.9,
    align=center,
    inner sep=1.0pt,
    minimum width=22mm,
    font=\zihao{-5},
    blur shadow={shadow blur steps=4, shadow blur extra rounding=1.0pt, shadow xshift=0.8pt, shadow yshift=-0.8pt, shadow opacity=16}
  },
  flowline/.style={-{Stealth[length=2.8mm,width=1.8mm]}, line width=0.65pt, draw=blue!70!black},
  flowlabel/.style={font=\zihao{-5}, text=blue!70!black, fill=white, inner sep=0.6pt}
}
\renewcommand*{\mkbibsuperscript}[1]{\textsuperscript{\zihao{-4}[#1]}}
\let\cite\supercite
\renewcommand*{\bibfont}{\songti\zihao{-4}\setstretch{1.5}}
\setlength{\bibhang}{2em}
\newcommand{\thesisfig}[2]{%
  \IfFileExists{#1}{%
    \includegraphics[width=#2\textwidth]{#1}%
  }{%
    \fbox{%
      \parbox[c][0.28\textheight][c]{#2\textwidth}{\centering
      待插入图片\\[4pt]\texttt{\detokenize{#1}}}%
    }%
  }%
}
\newcommand{\thesisfigfit}[3]{%
  \IfFileExists{#1}{%
    \includegraphics[width=#2\textwidth,height=#3\textheight,keepaspectratio]{#1}%
  }{%
    \fbox{%
      \parbox[c][#3\textheight][c]{#2\textwidth}{\centering
      待插入图片\\[4pt]\texttt{\detokenize{#1}}}%
    }%
  }%
}

% 表格内容字体统一：宋体五号（不影响图注/表注）
\AtBeginEnvironment{tabular}{\songti\zihao{5}}
\AtBeginEnvironment{longtable}{\songti\zihao{5}}

\makeatletter
\newcommand{\thesistableofcontents}{%
  \clearpage
  \phantomsection
  \pdfbookmark[0]{目录}{toc}%
  \begin{center}
  {\heiti\zihao{3}目录}
  \end{center}
  \@starttoc{toc}%
  \clearpage
}
\makeatother

\title{家用安防与智能家居联动系统设计}
\author{作者姓名}
\date{\today}

\begin{document}

\pagestyle{empty}
\IfFileExists{../规范与封面/论文封面.pdf}{%
  \includepdf[pages=1,pagecommand={},fitpaper=true]{../规范与封面/论文封面.pdf}%
}{%
  % 仓库未包含封面 PDF 时，跳过封面以保证可编译输出
}
\clearpage
\thispagestyle{empty}
\null
\clearpage

\frontmatter
\pagestyle{fancy}
\clearpage
% 封面不计入页码体系，前置部分从 I 开始连续编码
\setcounter{page}{1}
\pagenumbering{Roman}
\thesistableofcontents

\begin{center}
{\heiti\zihao{3}摘要}
\end{center}
\phantomsection
\addcontentsline{toc}{chapter}{摘要}
{\setlength{\parindent}{2em}
前期调研和样机预研中可以看到，家庭安防场景的关键问题不只是设备能不能联网。真正影响系统可用性的，是门禁、入侵和燃气应急三条安全链路在弱网条件下能不能稳定闭环。针对这一问题，本文设计并实现了一套以STM32F103RCT6为核心的单芯片家用安防与智能家居联动系统。系统能够完成密码/指纹双通道门禁验证，布防状态下的PIR入侵检测，MQ-2燃气异常判定，DHT11温湿度采集，继电器开锁与燃气阀切断，蜂鸣器本地告警，以及OV7670+FIFO触发抓拍并写入SD卡。网络侧利用ESP8266接入OneNET，通过Token鉴权的MQTT TCP链路完成温湿度、燃气、门锁、布防、告警等状态上报，并支持移动端远程开锁、布撤防和阈值配置。本地敏感数据由AES-128保护，App到云平台使用HTTPS/TLS加密，设备端MQTT TLS作为后续增强路径保留。

测试阶段按单元功能、系统联动、性能指标、稳定性与异常工况四层展开。实测结果为：门禁识别准确率98.0\%，入侵误报率3.2\%，本地告警平均时延620\,ms，远程控制平均时延1850\,ms，本地抓拍写入成功率96.5\%，系统连续运行72\,h未出现死机和任务阻塞。测试结果说明，在本科毕设的工程规模内，单芯片集中控制方案可以让关键安全动作先在本地完成，同时保留远程协同能力。该方案在实时性、可靠性、可复现性和成本之间取得了较合适的平衡，可为低成本家庭安防联动系统提供工程参考。
}

\vspace{1\baselineskip}
\noindent{\heiti\zihao{4}关键词：}\,\zihao{-4}智能家居；家庭安防；STM32F103RCT6；FreeRTOS；MQTT

\vspace{0.5\baselineskip}
\clearpage
\phantomsection
\addcontentsline{toc}{chapter}{英文摘要}
\begin{center}
{\bfseries\zihao{3}Abstract}
\end{center}
\begin{sloppypar}
This project targets three practical household security paths: access control, intrusion response, and gas emergency handling. I implemented a single-board home security linkage prototype based on STM32F103RCT6, so that sensing, decision making, actuation, and reporting can run on one controller with clear timing behavior.

The hardware integrates AS608 fingerprint sensing, PIR, temperature-humidity sensing, MQ-2, OV7670+FIFO, SD storage, and relay/buzzer actuators. The firmware is organized in two stages: bare-metal mode for interface bring-up and timing checks, and FreeRTOS mode for concurrent business tasks. ESP8266 over USART3 provides token-authenticated MQTT over TCP connectivity to OneNET for telemetry and command exchange, while HTTPS/TLS is used between the mobile app and cloud platform. Sensitive local data is protected with AES-128, and MQTT TLS on the device side is reserved as an upgrade path.

Validation is conducted by unit tests, scenario linkage tests, metric evaluation, and long-run stress tests. Measured results are: 98.0\% fingerprint accuracy, 3.2\% intrusion false alarm rate, 620\,ms average local alarm latency, 1850\,ms average remote control latency, 96.5\% local capture write success rate, and stable 72-hour operation without task lockup. These results indicate that a single-controller architecture can keep critical safety actions local and deterministic while still supporting cloud collaboration, which is suitable for undergraduate-scale engineering implementation.
\end{sloppypar}

\vspace{1\baselineskip}
\noindent{\bfseries\zihao{4}Key words:}\,Smart home, Home security, STM32F103RCT6, FreeRTOS, MQTT

\mainmatter
\pagestyle{fancy}
\setcounter{page}{1}

\chapter{绪论}
\section{研究背景与意义}
选题初期，本文先从家庭场景中真实存在的风险入手，进行文献调研和需求梳理。调研后可以看到两个趋势：一方面，家庭侧智能设备的使用比例持续上升；另一方面，居住场景中的安全事件仍然较多。CNNIC第54次报告显示，截至2024年6月，我国网民规模达到10.9967亿，互联网普及率为78.0\%，其中通过智能家居设备上网的用户占比为21.9\%\cite{cnnic54_2024}。国家消防救援局通报显示，2024年全国接报火灾90.8万起，其中居住场所火灾30.9万起，占建构筑物火灾的79\%\cite{china_fire_2024}。从这些数据看，本课题的重点并不只是设备是否联网，而是安全链路在弱网条件下能否保持稳定，并形成可验证的本地闭环。

在预研样机和分模块联调时，本文进一步找到了三个工程痛点。第一，单个模块可以独立跑通，但跨模块联动不够稳定。第二，网络波动会明显拉长远程控制时延，使用体验也会受到影响。第三，异常状态缺少统一的降级和恢复路径。针对这些问题，本文确定了“本地安全动作优先、远程协同增强”的实现思路，也就是把门禁、入侵告警和燃气应急等关键链路固定在本地执行，再通过MQTT完成状态上报和远程协同。

结合测控技术与仪器专业的培养目标，本课题研究意义主要体现在三个方面。第一，本文把传感检测、嵌入式控制、通信和系统集成放在同一工程对象中完成闭环验证。第二，本文用量化指标约束设计，并用实测数据验证方案有效性。第三，本文形成了在本科毕设资源条件下可复现、可扩展的低成本实现路径，可为同类家用安防联动项目提供参考。
\section{国内外研究现状}
从研究内容看，家庭安防与智能家居联动系统大致经历了“单点检测—联网控制—边缘协同”的发展过程。早期方案通常围绕单个传感器或单一执行器展开，例如用烟雾传感器判断燃气风险，用PIR检测人体移动，用门磁或键盘判断门禁状态。这类方案的优点是结构简单、成本低、调试直观，比较适合8位单片机或低端MCU平台；但系统能力也受到限制。单传感器阈值只能反映某一类风险，无法区分短时干扰、真实异常和设备故障。比如PIR容易受到安装角度、热源和电源波动影响，MQ-2需要预热且基线会随环境变化漂移，温湿度传感器也会因为单总线时序和线缆质量出现偶发读数失败。因此，单点检测系统在演示阶段容易跑通，但进入家庭场景后，还需要滤波、防抖、状态锁定和异常恢复机制，才能减少误报和漏报。

随着物联网技术发展，国外关于智能家居和家庭安防的研究逐渐把重点放到设备互联、云端管理和跨终端控制上。Al-Fuqaha等对物联网关键技术、通信协议和应用模式进行了系统梳理，指出感知层、网络层和应用层的分层设计是智能家居系统扩展的重要基础\cite{smart_home_survey_2017}。在这类框架下，家庭终端不再只是本地报警器，而是可以通过MQTT、HTTP等协议接入云平台，完成设备注册、状态发布、远程控制和用户提醒。Raut等围绕智能家居安防与家居自动化设计了IoT系统，说明联网后可以把门禁、照明和报警功能统一到移动端界面中，提升用户管理便利性\cite{home_security_iot_2020}。这类研究推动了智能家居从“设备能动作”向“设备可管理、状态可追踪”发展。

云平台主导型方案的优势比较明显。第一，平台可以统一管理多个终端，便于设备发现、数据存储、远程控制和告警推送。第二，移动端界面开发成本较低，用户可以通过手机查看设备状态，也可以远程执行开锁、布防和阈值设置等操作。第三，云端保留历史数据后，后续可以扩展异常统计、场景分析和多设备联动。近年的分层IoT架构研究和轻量化网关研究，也都把云端管理和边缘网关作为家庭安防系统的重要组成部分\cite{bozorgi2026_iot_arch,serepas2025_gateway}。不过，这类方案也带来一个实际问题：如果关键安全动作完全依赖云端确认，那么网络质量会直接影响系统安全性。门锁开关、燃气阀切断和蜂鸣告警这类动作，需要在现场快速完成。如果设备处于弱网、路由器重启或云平台连接抖动状态，远程链路可能出现秒级延迟，严重时还会失去控制。因此，云平台更适合承担状态展示、远程配置和辅助控制，不宜成为本地安全动作的唯一前提。

为了解决云端依赖过强的问题，国外研究中逐渐出现边缘控制增强的思路。边缘控制并不是完全取消云端，而是把实时性要求高、与安全直接相关的判断放到本地设备或边缘网关执行。这样，即使网络暂时不可用，设备仍可以完成基础安全动作，待网络恢复后再同步状态。生物特征门禁研究中，认证方式也从单一密码扩展到指纹、人脸等多模态认证。Alkanan等研究了面向智能家居的多模态生物特征门禁方法，强调身份识别应兼顾安全性和便捷性\cite{alkanan2026_biometric}。在环境安全方面，Yusuf等针对居民燃气泄漏监测与应急响应进行了IoT化设计，强调燃气异常不应只停留在数据上报，还要联动阀门、蜂鸣器和用户提示形成闭环\cite{yusuf2026_gas_leakage}。这些研究说明，智能安防系统的核心不只是“检测到异常”，还包括“异常后能否及时处置”。

从安全治理角度看，国外研究还比较重视IoT设备本身的安全风险。智能家居设备长期在线，且通常部署在家庭网络环境中，一旦设备认证、固件更新或通信加密处理不当，就可能带来隐私泄露和非法控制风险。相关IoT安全研究指出，弱口令、明文传输、固件漏洞和接口暴露都是消费级物联网设备常见问题\cite{mqtt_iot_review,rivas2026_vulnerability}。因此，比较完整的智能家居方案不仅要实现传感、控制和联网，还要考虑身份认证、数据保护和异常访问防护。对于资源受限的MCU设备来说，全部采用复杂安全协议会增加RAM、Flash和调试压力，但至少应保证设备接入具有Token或密钥校验，移动端到云平台使用HTTPS/TLS，本地敏感参数不以明文方式随意保存。这一点也为本文后续的AES-128存储、OneNET Token鉴权和TLS扩展预留提供了依据。

国内研究更多面向本科教学、工程应用和低成本落地，常见技术路线是以STM32为核心控制器，外接温湿度、烟雾、人体红外、继电器、蜂鸣器、显示屏和无线通信模块，再通过OneNET、阿里云或自建服务器完成远程监测与控制。连志强等总结了智能家居安防系统关键技术与发展趋势，认为多传感采集、无线通信、平台接入和移动端交互仍是低成本方案的重要组成部分\cite{lian2025_trend}。王晨等设计了基于STM32的家庭安防联动控制系统，将传感器检测、执行器动作和远程控制结合起来，体现出MCU集中控制在小型系统中的优势\cite{wang2024_stm32_linkage}。岳峰等、蔡文博等也分别围绕STM32与物联网平台、STM32F103智能家居控制系统开展了实现研究，说明这类方案在硬件成本、资料获取和教学验证方面具有较高可行性\cite{yue2022_stm32_iot_home,cai2023_stm32f103_control}。

在云平台应用方面，国内不少研究选择OneNET作为设备接入平台。OneNET提供物模型、属性上报、命令下发和数据可视化能力，适合把嵌入式设备快速接入手机端或Web端。陈志远等基于OneNET平台实现了智能安防系统，将设备状态上报和平台控制结合起来，为低成本设备接入云端提供了参考\cite{chen2023_onenet_ai}。这类方案的优势是开发周期短，平台文档较完整，适合毕业设计和中小型项目验证。但在实际使用中，OneNET等平台通常只能保证控制命令到达和属性记录，并不能替代设备端本地业务逻辑。比如平台返回命令接收成功，不等于门锁已经实际打开，也不等于本地告警动作已经完成。因此，设备端仍需要执行确认、状态回写和异常兜底，否则系统容易出现“平台显示成功、现场动作失败”的问题。

在无线通信和组网方面，国内研究也给出了不同实现路线。王磊等针对无线通信家庭安防系统进行了设计，说明WiFi、蓝牙、ZigBee等方式都可以用于家庭场景\cite{wang2024_home_security_wireless}。纪航等研究了基于STM32与ZigBee的家庭安防监测系统，体现出ZigBee在低功耗、多节点组网中的优势\cite{ji2024_stm32_zigbee}。不过，本课题的目标不是构建大规模多节点网络，而是在一套单芯片样机中完成门禁、入侵、燃气、环境采集和移动端控制闭环。与ZigBee多节点方案相比，ESP8266 WiFi接入更容易直接对接OneNET平台，也方便在实验室和宿舍网络环境中调试。它的不足是弱网和重连时序更明显，因此需要在固件中设置重连、超时、上报节流和本地降级策略。

国内关于环境监测与安防联动的研究还强调了传感器融合和现场标定。董航等面向家庭场景设计环境监测与安防联动系统，说明温湿度、气体检测和人体触发可以共同服务于家庭安全判断\cite{dong2024_env_monitor}。但从实际工程看，多传感器并不意味着简单把多个模块接到同一块板上。不同传感器的采样周期、噪声特性和触发方式差异较大：DHT11属于低速单总线数字传感器，适合周期采样；MQ-2属于模拟气敏器件，需要预热、滤波和阈值标定；PIR输出数字触发信号，但需要处理保持时间和误触发；摄像头和SD卡则会引入更高的瞬时带宽和供电负载。若没有统一的任务调度和事件管理，多个模块同时工作时容易出现采样阻塞、串口冲突、显示刷新混乱和告警动作延迟。

从已有文献和常见毕设实现看，当前家庭安防联动系统仍存在几个不足。第一，部分方案重功能演示，轻过程验证。论文中往往能看到传感器读数、手机端控制或继电器动作，但对测试过程、重复次数、失败样本和响应时间分布记录不够。第二，部分方案重联网，轻本地自治。系统可以远程点灯、开锁或查看数据，但网络断开后关键动作如何执行、恢复后状态如何同步，说明不够完整。第三，部分方案重模块堆叠，轻资源协调。STM32这类MCU资源有限，多串口、摄像头、SD卡、OLED和传感器同时接入后，任务优先级、缓冲区、临界区和看门狗都需要统一考虑。第四，部分方案对安全机制描述较少，本地密码、平台Token、移动端HTTPS和设备异常复位后的安全态需要形成闭环，而不是只在结论中笼统说明“系统安全可靠”。

根据上述研究现状，本文把课题重点放在“本地自治+远程协同”的落地验证上。系统使用STM32F103RCT6集中接入密码键盘、AS608指纹模块、HC-SR501人体红外、MQ-2燃气传感器、DHT11温湿度传感器、OV7670+FIFO摄像头、SD卡和继电器/蜂鸣器等模块。门禁验证、入侵告警、燃气应急和抓拍落盘优先在本地完成，保证弱网或断网时关键动作不被阻塞；OneNET和移动端主要承担状态展示、远程配置、控制确认和后续图片链路扩展。与单模块演示相比，本文更关注多模块并发后的稳定性、异常恢复和测试可复现性；与复杂网关方案相比，本文控制硬件规模和成本，使其更适合本科毕业设计条件下实现和验收。
\section{研究目标与主要内容}
本文的研究目标，是完成一套能够实际搭建、运行和验证的家庭安防联动系统。系统不把重点放在单个传感器演示上，而是面向家庭门厅、客厅和厨房等常见区域，验证门禁、入侵、燃气应急、环境监测和远程协同能否在同一块控制板上稳定配合。研究中坚持“关键动作本地优先，联网能力作为协同增强”的原则，使系统在网络波动时仍能完成基础安全处置。

围绕上述目标，本文主要完成四方面工作：第一，完成单芯片集中控制的总体方案设计，明确感知、判定、执行、通信和存储之间的关系；第二，完成门禁、入侵和燃气应急三条本地安全链路，使开锁、告警、关阀和状态记录能够在设备端闭环；第三，接入云平台和移动端，实现状态查看、远程控制和参数配置，但不让关键安全动作完全依赖网络；第四，按照工程验收过程开展功能、联动、性能和稳定性测试。具体功能拆分和量化指标放在第2章需求分析中集中说明，本节只给出研究目标和工作边界。

本文实施路线按实际开发过程推进。前期先完成需求拆分和模块选型，确定总体架构；随后完成硬件接线、电源分配、信号调理和接口资源核对；软件阶段先用裸机方式验证单模块时序，再切换到FreeRTOS组织多任务运行；最后通过单元测试、联动测试、性能测试和长稳测试检查系统是否达到第2章给出的需求。
\section{论文结构安排}
全文共7章，各章安排如下：第1章说明选题背景、文献调研结果、研究目标和技术路线；第2章给出系统需求分解、总体架构和关键技术选型；第3章介绍硬件实现、电气参数设计和调试过程；第4章说明软件架构、核心控制逻辑和通信安全机制；第5章介绍移动端交互设计与核心功能实现；第6章给出测试方案、实测数据和误差分析；第7章对课题工作进行总结，并提出后续优化方向。整体结构与“需求—设计—实现—验证”的工程流程保持一致。

论文整体章节结构如图\ref{fig:chapter_structure}所示。
\begin{figure}[H]
\centering
\begin{tikzpicture}[
  node distance=0.62cm,
  chapterbox/.style={
    rectangle,
    rounded corners=3pt,
    draw=black!55,
    fill=gray!8,
    minimum width=11.8cm,
    minimum height=0.72cm,
    align=center,
    font=\songti\zihao{5}
  },
  arrow/.style={-{Stealth[length=2.4mm]}, black!55, line width=0.6pt}
]
\node[chapterbox] (c1) {第1章：绪论——说明选题背景、研究现状与技术路线};
\node[chapterbox, below=of c1] (c2) {第2章：系统总体方案——完成需求分解、架构设计与技术选型};
\node[chapterbox, below=of c2] (c3) {第3章：硬件设计——给出模块电路、接口分配与联调排查};
\node[chapterbox, below=of c3] (c4) {第4章：软件设计——说明任务架构、控制逻辑与通信协议};
\node[chapterbox, below=of c4] (c5) {第5章：移动端与云端——实现状态查看、远程控制与抓拍入口};
\node[chapterbox, below=of c5] (c6) {第6章：系统测试——完成单元、联动、性能和稳定性验证};
\node[chapterbox, below=of c6] (c7) {第7章：总结与展望——归纳课题工作并提出后续优化方向};
\draw[arrow] (c1) -- (c2);
\draw[arrow] (c2) -- (c3);
\draw[arrow] (c3) -- (c4);
\draw[arrow] (c4) -- (c5);
\draw[arrow] (c5) -- (c6);
\draw[arrow] (c6) -- (c7);
\end{tikzpicture}
\caption{论文章节结构图}
\label{fig:chapter_structure}
\end{figure}

\chapter{系统总体方案设计}
\section{系统需求分析}
本章方案设计不是从理想智能家居平台出发，而是从当前已经搭建的实物系统出发。实物模型以一个封闭箱体模拟家庭局部空间：箱体前侧放置键盘、指纹模块、PIR和OLED，顶部和内部布置DHT11、MQ-2、摄像头和小型家居场景，底部连接主控板、电源、继电器、ESP8266和SD卡。这样的搭建方式决定了系统需求不能写得过大，重点应放在“这套样机能稳定完成哪些安全闭环”。

按实际使用区域划分，门厅位置主要对应密码/指纹开锁和异常尝试提示；客厅、过道位置主要对应PIR人体触发、蜂鸣器告警和摄像头抓拍；厨房或燃气区域主要对应MQ-2气体异常判断、燃气阀继电器切断和持续告警；移动端则承担远程开锁、布撤防、阈值配置和状态查看。由于家庭WiFi和云平台连接不一定稳定，本文在需求定义时把“弱网下关键动作仍可在本地完成”作为硬约束条件。也就是说，开锁、关阀、蜂鸣和本地抓拍不能等待云平台确认后才执行。
\subsection{功能需求}
结合前期搭建和联调过程，本文将系统功能归纳为四个可验证闭环。  
\kuan{1}门禁闭环：用户通过4$\times$4矩阵键盘输入密码，或通过AS608完成指纹识别。验证通过后，主控驱动门锁继电器进入开锁状态，同时OLED给出本地提示并更新门锁状态；验证失败时记录失败次数，达到阈值后进入告警处理。  
\kuan{2}入侵闭环：系统处于布防状态时，HC-SR501检测到人体移动后先形成候选事件，再由软件侧做持续时间和锁定窗口判断。事件成立后，系统触发蜂鸣器、本地日志、状态上报和OV7670+FIFO抓拍。当前验收阶段的图片闭环以SD卡生成BMP文件为准，不把图片上传到App作为通过条件。  
\kuan{3}燃气应急闭环：MQ-2采样值超过阈值并通过滤波判断后，系统优先驱动燃气阀继电器执行切断动作，同时蜂鸣器持续提示，并通过OLED和OneNET回传异常状态。该链路不依赖移动端确认，避免网络延迟影响应急动作。  
\kuan{4}远程协同闭环：ESP8266接入OneNET后，移动端可以查看温湿度、燃气、门锁、布防和告警状态，也可以下发开锁、布撤防和阈值配置命令。设备端收到命令后先在本地执行，再进行状态回传，避免只停留在“平台下发成功”的表面结果。
\subsection{性能需求}
承接第1章提出的研究目标，性能需求只保留可测试、可判定的量化指标。门禁链路看识别准确率，入侵链路看误报率和本地告警时延，远程链路看指令响应时延，整机运行看功耗和72\,h稳定性，图像链路在当前阶段看本地BMP写入是否成功。对应指标如下：  
\yanjiuitem{1}{门禁识别准确率$\geq95.0\%$；}
\yanjiuitem{2}{入侵检测误报率$\leq5.0\%$；}
\yanjiuitem{3}{本地告警响应时延$\leq1000$\,ms；}
\yanjiuitem{4}{远程控制响应时延$\leq2500$\,ms；}
\yanjiuitem{5}{空闲功耗$\leq600$\,mW，告警状态功耗$\leq1200$\,mW；}
\yanjiuitem{6}{本地抓拍写入成功率$\geq90.0\%$；}
\yanjiuitem{7}{连续运行72\,h无死机与任务阻塞。}
这些指标用于约束后续设计和第6章测试：门禁准确率对应身份验证链路；入侵误报率和告警时延对应PIR防抖、抓拍复核、任务优先级和执行器驱动；远程控制时延对应ESP8266重连、MQTT解析和回执机制；本地抓拍写入成功率对应OV7670+FIFO、SD卡和FatFs写入链路。这样处理后，第1章只说明“要解决什么问题”，本节负责把目标转化为可验收的指标。
\subsection{工程约束}
当前样机还受到几项具体工程约束。第一，主控使用STM32F103RCT6，GPIO、串口、SPI、ADC和外部中断资源需要同时分配给AS608、ESP8266、SD卡、OV7670、MQ-2、DHT11、PIR、键盘、OLED和执行器，不能随意增加新的控制板。第二，OV7670+FIFO、SD卡写入和ESP8266联网都会带来较明显的瞬时负载，摄像头长杜邦线、SD卡供电和继电器动作都可能影响稳定性。第三，AS608与串口调试存在引脚复用风险，系统调试时必须避免日志输出干扰指纹模块。第四，OneNET平台适合完成属性上报和命令下发，但图片上传链路受平台字段长度、网络带宽和设备端缓存限制，当前阶段采用本地SD抓拍优先的收尾方案。

因此，本文在方案设计中不追求“所有功能都上云”，而是按优先级取舍：门禁、燃气、入侵告警和本地抓拍属于本地必须完成的安全链路；状态展示、远程控制和参数配置属于联网增强链路；图片上传和App在线预览作为后续扩展保留。这样处理后，系统既能满足当前毕业设计验收，也不会因为网络图片链路不稳定而影响主功能判断。
\section{系统总体架构}
结合实物连线和软件任务划分，系统总体架构可以分为现场感知、主控决策、本地执行、本地存储、云端通信和移动端交互六部分。现场感知部分包括键盘、AS608、PIR、DHT11、MQ-2和OV7670+FIFO，分别提供身份输入、人体触发、环境状态、气体异常和图像帧数据。主控决策部分由STM32F103RCT6承担，负责把不同输入转化为门禁状态、布防状态、告警状态和抓拍请求。本地执行部分包括门锁继电器、燃气阀继电器、蜂鸣器和OLED，负责把判定结果直接反馈到现场。本地存储部分由SD卡和FatFs完成，主要保存抓拍BMP和日志信息。云端通信部分由ESP8266和OneNET MQTT链路承担，完成属性上报和命令下发。移动端负责状态展示、远程开锁、布撤防和阈值配置。

这套架构的关键点在于数据流和控制流分开处理。传感数据从现场模块进入STM32，先在本地完成快速判定；只有状态结果、告警标志和控制反馈才进入OneNET。移动端命令从云端下发后，也必须回到STM32本地状态机执行，而不是直接驱动外设。这样可以避免网络状态和本地执行状态不一致，也能保证断网时本地安全链路仍然运行。
系统总体结构如图\ref{fig:system_arch}所示。
\begin{figure}[H]
\centering
\thesisfig{images/fig2_1_system_arch.png}{0.96}
\caption{系统总体架构}
\label{fig:system_arch}
\end{figure}
\section{单芯片控制架构设计}
本课题没有采用“主控板+通信板+图像板”的多控制器方案，而是使用STM32F103RCT6完成集中控制。这样做的原因有三点。第一，当前系统外设数量虽然多，但大多数模块的数据量并不大，除摄像头抓拍外，键盘、指纹、PIR、DHT11、MQ-2和继电器都可以由单片机统一管理。第二，多控制器会增加板间通信协议和同步问题，反而拉长联调时间。第三，毕业设计验收更看重功能闭环和测试可复现性，单芯片方案更容易说明每条链路的状态来源和执行结果。

在资源分配上，STM32需要同时处理多类接口：AS608占用USART1，ESP8266使用USART3，SD卡走SPI2，OV7670+FIFO使用并口和控制信号，MQ-2占用ADC输入，PIR使用外部中断或轮询兜底，DHT11走单总线，OLED使用软件I2C。为了避免各模块互相抢占，软件上采用“安全任务优先、网络任务次之、显示和日志后台执行”的策略。门禁、燃气和入侵告警优先保证本地执行；网络上报和移动端命令解析可以延后；OLED刷新和日志整理不应阻塞关键动作。

图像链路是单芯片方案中最容易出问题的一环。OV7670直接输出图像数据对MCU时序要求较高，因此本系统使用带FIFO的摄像头模块，把摄像头采样和MCU读取解耦。当前验收口径下，抓拍后的图像先写入SD卡形成BMP文件，OLED显示文件大小作为写入成功提示。这样可以先证明“触发—采集—落盘”链路可用，再考虑HTTP图床或MQTT分包上传。
\section{系统工作流程}
本文将系统工作流程拆分为上电初始化、常态监测、事件触发、本地联动和远程同步五个阶段。这里的流程对应当前样机的实际运行顺序，而不是通用智能家居流程。

上电后，主控先完成时钟、GPIO、USART、ADC、外部中断、SD卡和OLED初始化，再启动门锁、传感器、网络、显示和日志相关任务。门禁链路优先进入可用状态，保证本地密码和指纹验证不依赖网络；随后网络任务连接ESP8266和OneNET，用于状态上报与远程控制。若网络连接失败，系统不停止本地功能，而是在OLED或状态变量中保留网络异常提示。

常态监测阶段，DHT11按周期读取温湿度，MQ-2持续采集ADC并进行滤波，PIR通过中断和轮询兜底捕获人体触发，键盘和AS608等待用户输入。事件触发后，系统按事件类型进入不同路径：合法身份通过后执行开锁并记录日志；布防状态下PIR触发后执行蜂鸣、抓拍落盘和告警上报；MQ-2超过阈值后优先关阀并持续告警；移动端命令到达后由设备端解析并进入同一套本地状态机。网络异常时，开锁、关阀、蜂鸣和SD抓拍仍在本地完成，待网络恢复后再同步状态。

系统工作流程如图\ref{fig:system_workflow}所示。
\begin{figure}[H]
\centering
\thesisfigfit{images/fig2_2_workflow.png}{0.92}{0.72}
\caption{系统工作流程}
\label{fig:system_workflow}
\end{figure}
\section{关键技术选型}
技术选型阶段，本文主要看三件事：当前硬件能否支撑、毕设周期内能否真正落地、测试时能否拿出证据。按照这三个判断标准，本文从控制器、调度机制、通信链路和安全机制四个维度进行对比，结果如表\ref{tab:selection_matrix}所示。结合当前收尾约束，图像链路使用“两阶段口径”：方案一先完成本地SD抓拍验收，方案二再扩展HTTP图床或MQTT分包上传。
\begin{table}[H]
\centering
\caption{关键技术选型对比表}
\label{tab:selection_matrix}
\renewcommand{\arraystretch}{1.15}
\begin{tabular}{>{\raggedright\arraybackslash}p{2.5cm} >{\raggedright\arraybackslash}p{3.6cm} >{\raggedright\arraybackslash}p{3.6cm} >{\raggedright\arraybackslash}p{3.1cm}}
\toprule
选型维度 & 备选方案A & 备选方案B & 本文采用方案 \\
\midrule
控制器平台 & 51/ESP32等：成本低，但并口采集与多外设并发受限 & STM32F103系列：外设资源均衡，开发资料成熟 & STM32F103RCT6 \\
调度机制 & 裸机轮询：实现简单，但并发与优先级管理弱 & FreeRTOS：任务化并发、可维护性高 & FreeRTOS \\
通信链路 & 仅HTTP轮询/直传：实现直观，但弱网稳定性与扩展性一般 & MQTT状态上报+图片分阶段链路：先本地SD验收，后续扩展HTTP图床或MQTT分包（自定义Topic非\texttt{\$sys}） & MQTT主链路+分阶段图片方案 \\
安全机制 & 单一链路防护：覆盖面有限 & 本地+链路双层防护：AES-128本地存储保护，设备侧Token鉴权接入MQTT TCP，App到云平台HTTPS/TLS，设备端TLS可扩展 & 双层防护方案 \\
\bottomrule
\end{tabular}
\end{table}
\clearpage
控制器选型时，本文重点核对GPIO资源、并口摄像头采集能力、多串口并发能力和定时中断负载，同时考虑开发资料、工具链和器件可获得性。STM32F103RCT6可以同时满足AS608串口、ESP8266串口、SD卡SPI、MQ-2 ADC、PIR输入、键盘扫描、OLED显示和OV7670+FIFO控制需求，因此作为当前样机的主控更合适\cite{stm32_ref,stm32_rm0008}。

在系统调度机制选型中，前期单模块调试使用裸机方式更直观，例如单独验证键盘、OLED、SD卡或摄像头时，可以减少任务切换带来的干扰。进入整机联调后，门锁、传感器、网络、显示和日志需要并发运行，裸机超循环难以保证优先级和超时管理。因此，最终运行版本使用FreeRTOS组织任务，同时保留裸机路径用于排障\cite{freertos_book,freertos_sched}。

通信链路需要同时考虑低带宽开销、弱网可用性和异步消息分发。与HTTP轮询相比，MQTT在报文开销、连接保持和发布/订阅解耦方面更适合作为主链路，所以本文使用OneNET MQTT承担温湿度、燃气、门锁、布防、告警等状态上报与控制下发。图片链路按工程节奏采用“先本地、后联网”的两阶段方案：先完成本地SD抓拍闭环验收，后续版本再扩展HTTP图床或MQTT分包上传。其中，MQTT分包使用平台允许的自定义Topic，不使用\texttt{\$sys}保留命名空间\cite{mqtt_spec,mqtt_v5}。

在安全机制选型中，安全方案需要同时覆盖本地存储和公网传输。考虑MCU资源约束，本文使用AES-128保护本地敏感数据；App到云平台的HTTP接口使用HTTPS/TLS加密，设备端则通过Token鉴权接入OneNET MQTT TCP链路，并将MQTT TLS作为后续增强方案，形成“本地数据保护+网络链路可扩展防护”的双层设计\cite{aes_fips197,rfc5246,tls_rfc}。

\section{本章小结}
本章从当前实物样机出发，完成了系统需求拆分、总体架构设计、单芯片控制策略和关键技术选型。方案设计中明确了三条主线：第一，门禁、入侵和燃气应急必须在本地闭环；第二，OneNET和移动端用于远程协同，不替代本地安全动作；第三，图像链路先以SD卡BMP落盘作为验收依据，上传展示作为后续增强。上述方案为第3章硬件连接、第4章软件状态机和第6章测试验证提供了统一设计基线。

\chapter{系统硬件设计}
\section{硬件总体设计}
硬件设计阶段，本文先明确三项工程目标：关键链路能够闭环，接口边界便于调试，关键参数可以计算。按照这些目标，系统以STM32F103RCT6为控制核心，把感知、执行、通信和存储单元集中接入。结构上采用单板集中控制，减少跨板连接和协议同步带来的联调风险；电气上采用分域供电和驱动隔离，重点抑制继电器动作、无线发射对采样和控制链路的扰动。最终方案满足测控系统“测量稳定、控制可靠、异常可恢复”的基础要求。
硬件总体连接关系如图\ref{fig:hw_block}所示。

\begin{figure}[H]
\centering
\thesisfig{images/fig3_1_hw_block.png}{0.85}
\caption{硬件总体连接框图}
\label{fig:hw_block}
\end{figure}

\begin{table}[htbp]
\centering
\caption{系统硬件模块组成与功能}
\renewcommand{\arraystretch}{1.15}
\begin{tabular}{>{\raggedright\arraybackslash}p{3.0cm} >{\raggedright\arraybackslash}p{4.1cm} >{\raggedright\arraybackslash}p{3.1cm} >{\raggedright\arraybackslash}p{2.8cm}}
\toprule
模块类别 & 核心器件 & 主要功能 & 备注 \\
\midrule
主控模块 & STM32F103RCT6（单芯片） & 任务调度、网络通信、联动决策 & 裸机/FreeRTOS可切换 \\
门禁模块 & AS608 & 指纹采集与身份验证 & UART \\
安防检测模块 & HC-SR501 & 人体移动检测 & GPIO中断 \\
环境监测模块 & DHT11、MQ-2 & 温湿度/气体采集 & 周期采样 \\
图像采集模块 & OV7670+FIFO & 异常抓拍与图像缓存 & 并口 \\
执行模块 & 门锁/燃气阀继电器、蜂鸣器 & 设备联动与本地告警 & 驱动隔离 \\
存储模块 & SD卡（SPI2） & 抓拍BMP与容量统计 & FatFs \\
\bottomrule
\end{tabular}
\end{table}
\section{主控最小系统设计}
主控最小系统部分，本文围绕“稳定启动、稳定运行、便于调试”展开。系统由时钟、复位、调试下载和电源去耦四个基础电路构成，目标是在72\,MHz主频下保持长期稳定并具备抗瞬态扰动能力。时钟使用外部晶振+PLL来满足串口通信与定时采样精度；复位使用上电复位与按键复位并行的方式，避免慢上电导致的启动异常；调试接口使用SWD（PA13/PA14），在保留在线调试能力的同时释放JTAG引脚资源供业务功能使用。
主控最小系统原理结构如图\ref{fig:mcu_min_system}所示。

\begin{figure}[H]
\centering
\thesisfig{images/fig3_2_mcu_min_system.png}{0.78}
\caption{主控最小系统原理图}
\label{fig:mcu_min_system}
\end{figure}

\noindent 如图\ref{fig:sch_stm32_min}所示，为便于复现硬件连接细节，本文同时给出主控最小系统的完整原理图。  
\begin{figure}[H]
\centering
\thesisfigfit{images/原理图/sch_stm32f103rct6.png}{0.96}{0.62}
\caption{STM32F103RCT6主控最小系统原理图（原理图）}
\label{fig:sch_stm32_min}
\end{figure}

\Needspace{10\baselineskip}
\noindent 主控最小系统的关键配置与接口资源占用如表\ref{tab:mcu_min_params}所示。

\begin{table}[htbp]
\centering
\caption{主控最小系统关键参数}
\label{tab:mcu_min_params}
\renewcommand{\arraystretch}{1.15}
\begin{tabular}{>{\raggedright\arraybackslash}p{3.3cm} >{\raggedright\arraybackslash}p{2.9cm} >{\raggedright\arraybackslash}p{6.4cm}}
\toprule
参数项 & 配置 & 说明 \\
\midrule
系统时钟 & 72MHz & 外部晶振+PLL \\
工作电压 & 3.3V & 稳压供电 \\
下载调试接口 & SWD（PA13/PA14） & 关闭JTAG，保留SWD \\
主要串口 & USART1 + USART3 & USART1用于AS608/调试，USART3用于ESP8266 \\
核心外设占用 & OV并口+SPI2+多GPIO & 支持摄像头、SD、传感与执行器并发 \\
\bottomrule
\end{tabular}
\end{table}
\subsection{主控去耦与时钟稳定设计}
针对主控高速切换电流，电源引脚按“就近、小容值高频旁路+中容值储能”配置去耦网络：每组VDD/VSS邻近放置0.1\,$\mu$F陶瓷电容，并在电源入口配置10\,$\mu$F电解/钽电容。其设计依据为
\begin{equation}
C \ge \frac{I_{\mathrm{step}}\cdot \Delta t}{\Delta V}
\label{eq:decouple}
\end{equation}
其中，$I_{\mathrm{step}}$为瞬态电流，$\Delta t$为电流跃迁持续时间，$\Delta V$为允许电压跌落。按$I_{\mathrm{step}}=60$\,mA、$\Delta t=1$\,$\mu$s、$\Delta V=0.1$\,V估算，最小去耦电容约为0.6\,$\mu$F。若采用$n$颗0.1\,$\mu$F电容并联，其总容值为
\begin{equation}
C_{\mathrm{sum}}=n\times 0.1\ \mu \mathrm{F}
\label{eq:decouple_parallel}
\end{equation}
工程上取$n=6$时$C_{\mathrm{sum}}=0.6\ \mu\mathrm{F}$，与式\eqref{eq:decouple}估算下限一致，并叠加10\,$\mu$F储能电容提升低频负载扰动抑制能力。时钟支路通过短走线、对称布线和地参考完整性控制降低时钟抖动与串扰。
\noindent 去耦电容与电源入口储能电容的连线关系如图\ref{fig:mcu_decouple_network}所示。
\begin{figure}[H]
\centering
\thesisfigfit{images/fig3_2_1_mcu_decouple.png}{0.92}{0.52}
\caption{主控电源去耦网络连线示意图}
\label{fig:mcu_decouple_network}
\end{figure}
\section{感知层模块设计}
感知层实现中，本文把“输入数据可信度”作为首要约束。按信号特性，本文将传感模块分为数字触发链路（AS608、PIR、温湿度）和模拟采样链路（MQ-2），先分别完成时序与电气稳定，再在控制层执行时间同步、阈值判定和事件融合，最终通过串口JSON或MQTT统一上报\cite{as608_manual,hcsr501_datasheet,mq2_datasheet,ov7670_datasheet}。这种分层处理方式便于在调试时快速判断问题来源于采集端还是判定端。
感知层各传感器接线关系如图\ref{fig:sensor_wiring}所示。

\begin{figure}[H]
\centering
\thesisfig{images/fig3_3_sensor_wiring.png}{0.78}
\caption{传感器模块接线示意}
\label{fig:sensor_wiring}
\end{figure}
\Needspace{6\baselineskip}
\subsection{AS608指纹模块}
在门禁模块选型中，本文重点比较了识别准确率、模板容量、接口实现复杂度和供电兼容性。AS608支持UART通信与模板管理，驱动复杂度低于并行接口方案，且与3.3\,V控制域兼容，满足本课题“低改造成本+快速联调”的实现需求。结合毕设阶段的功能规模，AS608相较仅支持小容量模板的低端模块更适合本系统。
\noindent AS608模块的实物形态与接线关系如图\ref{fig:photo_as608}和图\ref{fig:sch_as608_arch}所示。  
\begin{figure}[H]
\centering
\thesisfig{images/实物图/photo_as608.png}{0.52}
\caption{AS608指纹模块实物图（实物图）}
\label{fig:photo_as608}
\end{figure}
\begin{figure}[H]
\centering
\begin{tikzpicture}[
  node distance=18mm,
  block/.style={rounded corners=3pt, draw=blue!60!black, fill=blue!6, line width=0.8pt, minimum width=35mm, minimum height=13mm, align=center, font=\zihao{-4}},
  pin/.style={font=\zihao{5}, align=left},
  arrow/.style={-{Stealth[length=2.4mm]}, thick, draw=blue!65!black}
]
\node[block] (mcu) {STM32F103RCT6\\USART1 + GPIO};
\node[block, right=32mm of mcu] (fp) {AS608指纹模块\\模板采集/比对};
\draw[arrow] ([yshift=4mm]fp.west) -- node[above,pin]{TXD $\rightarrow$ PA10} ([yshift=4mm]mcu.east);
\draw[arrow] ([yshift=0mm]mcu.east) -- node[above,pin]{PA9 $\rightarrow$ RXD} ([yshift=0mm]fp.west);
\draw[arrow] ([yshift=-4mm]fp.west) -- node[below,pin]{PS\_STA $\rightarrow$ PA8} ([yshift=-4mm]mcu.east);
\node[pin, below=8mm of fp] (pwr) {VCC接3.3\,V，GND共地；\\通信波特率按工程配置为57600\,bps。};
\draw[blue!45!black, thick] (pwr.north) -- (fp.south);
\end{tikzpicture}
\caption{AS608模块接线与通信关系图}
\label{fig:sch_as608_arch}
\end{figure}
\Needspace{6\baselineskip}
\subsection{HC-SR501人体红外}
PIR模块部分，本文主要关注触发稳定性、部署便利性和误触发可调能力。HC-SR501输出数字电平，可直接接入EXTI中断链路，硬件复杂度低于外接放大与滤波的模拟方案；同时具备延时保持与灵敏度调节能力，便于按走廊/门厅场景进行参数收敛并抑制误报。
\noindent 从图\ref{fig:photo_hcsr501}可以看到，模块板载双可调电位器与触发模式跳帽：两路电位器分别用于灵敏度和延时参数整定，跳帽用于单次触发/可重复触发模式切换。工程实现中将其与软件侧布防状态机联合配置，可在不改硬件的前提下完成“响应速度-误报率”的现场折中。
\noindent HC-SR501模块实物与电路连接如图\ref{fig:photo_hcsr501}和图\ref{fig:sch_hcsr501}所示。  
\begin{figure}[H]
\centering
\thesisfig{images/实物图/photo_hcsr501.png}{0.55}
\caption{HC-SR501人体红外模块实物图（实物图）}
\label{fig:photo_hcsr501}
\end{figure}
\begin{figure}[H]
\centering
\thesisfigfit{images/原理图/sch_hcsr501.png}{0.82}{0.48}
\caption{HC-SR501人体红外模块原理图（原理图）}
\label{fig:sch_hcsr501}
\end{figure}
\Needspace{6\baselineskip}
\subsection{DHT11温湿度传感链路}
DHT11温湿度链路主要用于环境状态展示与联动阈值辅助判断。系统使用单总线温湿度采样方案，减少了模拟调理电路复杂度，能够满足家用场景下的周期监测需求。模块实物和单总线原理连接如图\ref{fig:photo_dht11}和图\ref{fig:sch_dht11_bus}所示。模块排针接口便于与STM32控制板快速连接，板载指示灯也方便联调时观察供电状态；DATA线通过4.7\,k$\Omega$上拉电阻接至VCC，保证总线空闲时保持高电平。联调阶段重点验证采样稳定性、超时恢复和异常重试策略，保证DHT11链路在多任务并发场景下可持续工作。
\begin{figure}[H]
\centering
\thesisfig{images/实物图/photo_dht11_real.png}{0.52}
\caption{DHT11温湿度传感器模块实物图（实物图）}
\label{fig:photo_dht11}
\end{figure}
\begin{figure}[H]
\centering
\thesisfig{images/原理图/sch_dht11_bus.png}{0.52}
\caption{DHT11单总线连接原理图（原理图）}
\label{fig:sch_dht11_bus}
\end{figure}
\Needspace{6\baselineskip}
\subsection{MQ-2烟雾}
MQ-2模块用于可燃气体与烟雾异常检测，选型依据为可获得性高、响应快且支持模拟量连续采样。与仅输出开关量的简化烟雾模块相比，MQ-2可提供ADC趋势数据，便于实现“阈值触发+变化率辅助判定”的双判据逻辑。
\noindent 从图\ref{fig:photo_mq2}可以看到，模块使用“敏感探头+小型接口板”一体化封装：上方金属防护网对应气敏探头，底部排针接口便于快速接入控制板并开展台架标定。该结构有利于在毕设阶段快速完成布线、预热、采样和参数整定闭环。
\noindent MQ-2传感器模块实物如图\ref{fig:photo_mq2}所示。  
\begin{figure}[H]
\centering
\thesisfig{images/实物图/photo_mq2.png}{0.62}
\caption{MQ-2烟雾/可燃气体模块实物图（实物图）}
\label{fig:photo_mq2}
\end{figure}

在信号调理方面，MQ-2输出经分压后接入ADC1\_IN8（PB0）。分压关系为
\begin{equation}
V_{\mathrm{ADC}} = V_{\mathrm{in}}\cdot \frac{R_2}{R_1+R_2}
\label{eq:mq2_divider}
\end{equation}
为使$V_{\mathrm{ADC}} \le 3.3$\,V，在传感器侧最高输出按5\,V估算时取$R_1=10$\,k$\Omega$、$R_2=20$\,k$\Omega$，则$V_{\mathrm{ADC,max}}\approx3.33$\,V。工程实现中进一步通过ADC满量程校准和软件限幅避免过压采样。模拟支路增加RC低通（典型$R=1$\,k$\Omega$，$C=0.1$\,$\mu$F）抑制高频噪声，其截止频率
\begin{equation}
f_c=\frac{1}{2\pi RC}\approx1.59\ \text{kHz}
\label{eq:rc_cutoff}
\end{equation}
满足低速环境量采样需求。其幅频响应为
\begin{equation}
\left|H(j\omega)\right|=\frac{1}{\sqrt{1+\left(\frac{f}{f_c}\right)^2}}
\label{eq:rc_mag}
\end{equation}
据式\eqref{eq:rc_mag}，当$f=20$\,kHz时$\left|H\right|\approx0.079$（约$-22.0$\,dB），当$f=50$\,kHz时$\left|H\right|\approx0.032$（约$-29.9$\,dB），说明该RC网络对高频尖峰噪声具有有效衰减能力。分压网络与RC联合后，ADC输入端纹波峰峰值在台架测试中由约120\,mV降至约34\,mV，降噪幅度约71.7\%。
\subsubsection*{气体传感器标定实验}
MQ-2烟雾/可燃气体传感器使用STM32F103RCT6内置12位ADC采集模拟量，系统供电5.0\,V，负载电阻$R_L=10$\,k$\Omega$，传感器预热30\,min至稳定后进行标定\cite{mq2_calib_source}。标定点原始数据如表\ref{tab:mq2_calib_raw}所示。

\begin{table}[htbp]
\centering
\caption{MQ-2传感器标定点原始数据}
\label{tab:mq2_calib_raw}
\renewcommand{\arraystretch}{1.15}
\begin{tabular}{>{\centering\arraybackslash}p{1.2cm} >{\centering\arraybackslash}p{2.4cm} >{\centering\arraybackslash}p{2.0cm} >{\centering\arraybackslash}p{2.2cm} >{\centering\arraybackslash}p{2.2cm}}
\toprule
等级 & 含义 & 参考浓度(ppm) & MQ-2输出电压(V) & STM32 ADC读数 \\
\midrule
0 & 洁净空气 & 0 & 0.18$\sim$0.25 & 145$\sim$205 \\
1 & 轻微油烟/泄漏 & $\sim$300 & 0.45$\sim$0.65 & 365$\sim$530 \\
2 & 明显烟雾/燃气 & $\sim$1000 & 0.90$\sim$1.30 & 735$\sim$1060 \\
3 & 高浓度告警 & $\sim$2000 & 1.60$\sim$2.10 & 1310$\sim$1720 \\
4 & 危险浓度 & $\geq$3000 & $\geq$2.40 & $\geq$1960 \\
\bottomrule
\end{tabular}
\end{table}

ADC值到传感器输出电压的转换公式为
\begin{equation}
V_{AO}=\frac{ADC}{4095}\times 5.0
\label{eq:mq2_adc_to_v}
\end{equation}
传感器等效电阻计算公式为
\begin{equation}
R_S = R_L \times \frac{V_{CC}-V_{AO}}{V_{AO}}
\label{eq:mq2_rs}
\end{equation}
其中$R_L=10$\,k$\Omega$，$V_{CC}=5.0$\,V。温漂补偿使用25$^\circ$C基准的线性修正：
\begin{equation}
A_{\mathrm{comp}} = A_{\mathrm{raw}}-\alpha_T(T-25),\quad \alpha_T=3.2\ \mathrm{ADC}/^\circ C
\label{eq:mq2_temp_comp}
\end{equation}
实测在洁净空气基线下，15$^\circ$C/25$^\circ$C/35$^\circ$C对应原始ADC均值约为172/205/239。经式\eqref{eq:mq2_temp_comp}补偿后，三组数据收敛至约204/205/207。
阈值设定使用“基线均值+3$\sigma$裕量”原则，即$T_{gas}=A_{\mathrm{base}}+3\sigma+\Delta_{\mathrm{safe}}$。本系统取$A_{\mathrm{base}}=205$、$\sigma=18$、$\Delta_{\mathrm{safe}}=40$，得到$T_{gas}=299$，工程实现中取300。通过标定可建立ADC读数与气体浓度等级映射关系，实现家用场景下烟雾/可燃气体分级检测与告警。
\Needspace{6\baselineskip}
\subsection{OV7670摄像头+FIFO}
图像采集链路选型重点是接口匹配性、带宽压力和抓拍确定性。OV7670本体数据吞吐较高，若直接由MCU实时搬运将显著占用总线与CPU时间；因此本文使用OV7670+FIFO方案，通过硬件缓存解耦采集与读取时序，保障告警任务优先执行\cite{ov7670_datasheet}。当前收尾阶段以本地SD抓拍验收为主，系统在按键或抓拍指令触发后将图像帧转换为RGB565 BMP并写入SD卡。HTTP图床和MQTT分包上传作为后续增强路径保留，不作为当前主线验收条件。
\noindent OV7670与AL422B FIFO的原理连接关系如图\ref{fig:photo_ov7670}所示。该连接通过FIFO缓存对摄像头并口数据进行时序解耦，降低MCU在抓拍瞬间的总线占用与中断抖动风险。  
\begin{figure}[H]
\centering
\thesisfigfit{images/原理图/sch_ov7670_fifo_al422b.png}{0.98}{0.60}
\caption{OV7670+FIFO图像采集链路原理连接图（原理图）}
\label{fig:photo_ov7670}
\end{figure}
\section{执行层驱动设计}
执行层负责将控制策略转化为物理动作，实现“检测-控制-反馈”闭环。执行单元驱动设计强调三点：电气隔离、反灌防护和可重复动作可靠性。
执行器驱动电路关系如图\ref{fig:actuator_driver}所示。

\begin{figure}[H]
\centering
\thesisfig{images/fig3_4_actuator_driver.png}{0.78}
\caption{执行器驱动电路示意}
\label{fig:actuator_driver}
\end{figure}
\subsection{继电器（门锁/燃气）}
继电器驱动模块需满足低压控制高压负载、触点容量匹配和反向电动势抑制。项目实现包含门锁继电器（PC6）与燃气阀继电器（PC7），照明联动位可复用PC8。驱动电路使用NPN三极管低端开关+续流二极管+可选光耦隔离。
\noindent 从图\ref{fig:photo_relay}可以看到，模块端子区明确标注了NO/COM/NC三触点，便于按“常开驱动”或“常闭保护”两种策略接入负载；板载状态指示灯可用于联调阶段快速确认吸合状态。继电器外壳标注的触点额定参数可作为负载选型与安全裕量校核依据。
\noindent 继电器执行模块实物如图\ref{fig:photo_relay}所示。  
\begin{figure}[H]
\centering
\thesisfig{images/实物图/photo_relay.png}{0.60}
\caption{继电器执行模块实物图（实物图）}
\label{fig:photo_relay}
\end{figure}

以5\,V继电器线圈电流$ I_C=70 $\,mA为例，三极管强迫放大倍数取$\beta_f=10$，则基极电流
\begin{equation}
I_B \ge \frac{I_C}{\beta_f}=7\ \text{mA}
\label{eq:ib}
\end{equation}
在GPIO输出3.3\,V、$V_{BE}\approx0.7$\,V条件下，基极电阻取
\begin{equation}
R_B = \frac{V_{GPIO}-V_{BE}}{I_B} \approx \frac{3.3-0.7}{7\ \text{mA}} \approx 371\ \Omega
\label{eq:rb}
\end{equation}
工程上选用390\,$\Omega$标准阻值，兼顾饱和裕量与GPIO驱动能力。续流二极管选用1N4148/1N4007并联在线圈两端，抑制关断尖峰，降低对MCU地参考的冲击。
\subsection{蜂鸣器告警}
蜂鸣器模块需满足3.3\,V兼容、声压可感知和节奏可编程要求。系统使用有源蜂鸣器低电平触发方案，通过定时器调制鸣叫占空比，实现入侵、燃气和门禁异常的差异化提示，避免单一长鸣导致用户判别困难。
\noindent 从图\ref{fig:photo_buzzer}可以看到，模块接口清晰标注了GND、I/O和VCC三引脚，且板面标识为低电平触发，便于与MCU控制逻辑直接对应。工程实现中将I/O口默认置高电平，仅在告警窗口按节奏拉低，从而兼顾静默待机与明确声提示。
\noindent 蜂鸣器驱动电路与模块实物分别如图\ref{fig:sch_buzzer_driver}和图\ref{fig:photo_buzzer}所示。  
\begin{figure}[H]
\centering
\thesisfig{images/原理图/sch_buzzer_driver.png}{0.62}
\caption{蜂鸣器驱动原理图（原理图）}
\label{fig:sch_buzzer_driver}
\end{figure}
\begin{figure}[H]
\centering
\thesisfig{images/实物图/photo_buzzer.png}{0.55}
\caption{蜂鸣器模块实物图（实物图）}
\label{fig:photo_buzzer}
\end{figure}
\section{板内接口与串口通信设计}
系统不使用主从板间串口协议，而是使用“板内多外设接口 + 外联串口通信”方案。核心串口分工为：USART1用于AS608及调试打印，USART3用于ESP8266联网，二者在软硬件上解耦，提升联调稳定性。
串口与外设数据流向如图\ref{fig:uart_dataflow}所示。

\begin{figure}[H]
\centering
\thesisfig{images/fig3_6_uart_dataflow.png}{0.99}
\caption{串口与外设数据流向图}
\label{fig:uart_dataflow}
\end{figure}

\subsection{ESP-01S（ESP8266）WiFi联网模块}
\noindent ESP-01S用于将设备接入家庭WiFi并对接OneNET MQTT服务，承担状态上报、告警推送与远程控制指令下发等通信任务。模块实物与原理连接如图\ref{fig:photo_esp01s}和图\ref{fig:sch_esp01s}所示。  
\noindent 从图\ref{fig:photo_esp01s}可以看到，模块右侧为板载PCB天线，左侧为引脚排针接口，便于在小体积条件下完成无线收发与串口接入。该结构适合本系统“主控板+外接通信模块”的分体式联调方式，既降低布板复杂度，也方便后期更换模块或排查链路问题。
\begin{figure}[H]
\centering
\thesisfig{images/实物图/photo_esp01s.png}{0.55}
\caption{ESP-01S（ESP8266）联网模块实物图（实物图）}
\label{fig:photo_esp01s}
\end{figure}
\begin{figure}[H]
\centering
\thesisfigfit{images/原理图/sch_esp01s.png}{0.92}{0.58}
\caption{ESP-01S（ESP8266）模块原理图（原理图）}
\label{fig:sch_esp01s}
\end{figure}
结合图\ref{fig:photo_esp01s}与图\ref{fig:sch_esp01s}，本系统ESP-01S联网模块使用USART3与主控通信，连线关系为PB10$\rightarrow$ESP\_RX、PB11$\leftarrow$ESP\_TX，模块使能脚CH\_PD/EN由PC9控制。默认工作波特率为115200\,bps，联网方式为STA模式，SSID与口令由统一宏定义管理，便于后期更换网络环境时快速调整。  

硬件启动约束上，ESP-01S需满足EN高电平与稳定3.3\,V供电后再进入AT初始化流程；软件侧先执行基础AT连通性检测，再进行入网与MQTT会话建立。为提高联调鲁棒性，工程保留了“预配置网络/固件直连配置”两种模式切换，并在启动阶段输出关键AT状态信息（如固件版本与连接状态），用于定位串口接线、波特率或入网失败等问题。

\subsection{SD卡存储模块}
\noindent SD卡用于抓拍BMP图片与运行日志的本地持久化存储，便于离线取证、断网场景下的证据保全与容量统计。当摄像头抓拍触发时，抓拍任务将FIFO读出的图像帧编码为BMP并落盘保存。当前默认方案不依赖图床或OneNET图片字段，先以SD根目录生成BMP文件和OLED显示文件大小作为验收依据；后续若开启上传，再从SD取图并走HTTP图床或MQTT分包路径。模块实物与原理连接如图\ref{fig:photo_sd}和图\ref{fig:sch_microsd}所示。  
\noindent 从图\ref{fig:photo_sd}可以看到，模块集成了MicroSD卡座与板载电平/接口电路，左侧排针便于直接接入MCU的SPI总线。该类模块化封装能够在不增加硬件设计复杂度的前提下快速搭建本地存储链路，满足抓拍文件落盘与离线取证需求。  
\begin{figure}[H]
\centering
\thesisfig{images/实物图/photo_sd.png}{0.55}
\caption{SD卡模块实物图（实物图）}
\label{fig:photo_sd}
\end{figure}
\begin{figure}[H]
\centering
\thesisfigfit{images/原理图/sch_microsd.png}{0.92}{0.58}
\caption{MicroSD读写模块原理图（原理图）}
\label{fig:sch_microsd}
\end{figure}
图\ref{fig:photo_sd}与图\ref{fig:sch_microsd}分别给出SD卡模块实物与原理连接。SD卡存储链路使用SPI2总线实现，SCK/MISO/MOSI分别为PB13/PB14/PB15，片选CS为PB12。文件系统使用FatFs，功能受\texttt{EN\_FATFS\_SD}编译开关控制：全功能模式默认为1；调试模式下仅在\texttt{DEBUG\_PASSWORD\_ONLY=1}时保留SD功能。为提高差卡与长杜邦线场景兼容性，初始化流程使用“上电稳定延时+80字节空时钟+ACMD41轮询重试”的保守参数，挂载稳定后再进入抓拍BMP落盘与容量统计流程。

\begin{table}[htbp]
\centering
\caption{关键硬件接口分配与电气说明}
{\zihao{5}\renewcommand{\arraystretch}{1.00}\begingroup\setlength{\tabcolsep}{3.5pt}%
\begin{tabular}{@{}>{\raggedright\arraybackslash}p{1.50cm} >{\raggedright\arraybackslash}p{2.68cm} >{\raggedright\arraybackslash}p{1.45cm} >{\raggedright\arraybackslash}p{2.72cm} >{\raggedright\arraybackslash}p{\dimexpr\textwidth-1.50cm-2.68cm-1.45cm-2.72cm-8\tabcolsep\relax}@{}}
\toprule
功能模块 & MCU引脚 & I/O类型 & 电气特性 & 设计说明 \\
\midrule
AS608指纹 & \makecell{PA9/PA10\\(USART1)} & UART & \makecell{3.3\,V TTL\\57600\,bps} & 指纹采集、模板比对与串口调试复用 \\
ESP8266联网 & \makecell{PB10/PB11\\(USART3)} & UART & \makecell{3.3\,V TTL\\115200\,bps} & MQTT上报与下行控制解耦 \\
PIR红外 & PC4 & EXTI & 数字输入，上拉，边沿触发 & 入侵事件中断触发，降低轮询负担 \\
DHT11 & PC5 & GPIO & 单总线，4.7\,k$\Omega$上拉 & 温湿度周期采样，时序软件校验 \\
MQ-2 & \makecell{PB0\\(ADC1\_IN8)} & ADC & \makecell{0--3.3\,V输入，\\12位ADC} & 分压+RC滤波后采样气体变化趋势 \\
SD卡（SPI2） & \makecell{PB13/PB14/PB15\\PB12(CS)} & SPI & \makecell{3.3\,V SPI\\512B扇区} & FatFs挂载、BMP抓拍存储与容量统计 \\
门锁/燃气阀 & PC6/PC7 & GPIO输出 & 三极管驱动，继电器隔离 & 本地执行链路，支持失网应急动作 \\
蜂鸣器 & PC10 & \makecell{GPIO\\PWM} & 3.3\,V驱动，限流保护 & 告警分级声提示与静音恢复 \\
\bottomrule
\end{tabular}\endgroup
}
\end{table}
\section{电源管理系统设计}
本系统实际使用DCP3512四路电源模块作为集中供电单元，输入为VIN 6--30\,V，板载电源开关控制总上电。模块提供四路输出：3.3\,V、5\,V、12\,V及一路可调输出（约0.9--29.7\,V），其中3.3\,V为STM32与数字传感链路供电，5\,V为部分外设与通信模块供电，12\,V用于门锁/阀体等执行器支路，可调路用于实验阶段外设兼容与电压验证。  

为减少执行器负载变化对采样精度影响，电源分配使用“数字/模拟/执行器分支”布线策略：ADC相关传感器优先接入3.3\,V低噪声分支，继电器与门锁等冲击负载独立走线并就近回流，降低公共回路压降耦合。关键采样点维持“0.1\,$\mu$F+1\,$\mu$F”局部去耦，执行器电源端增设47\,$\mu$F储能电容以抑制瞬态跌落。  

工程调试中通过逐路上电与分支电压复核（3.3\,V/5\,V/12\,V）确认模块输出稳定，再进行整机联调；当可调路参与实验供电时，先空载设定目标电压再接入负载，避免误调导致器件过压风险。
\begin{figure}[H]
\centering
\thesisfig{images/实物图/photo_power.png}{0.62}
\caption{电源模块实物图（实物图）}
\label{fig:photo_power}
\end{figure}
\begin{figure}[H]
\centering
\thesisfigfit{images/原理图/sch_power_module.png}{0.92}{0.60}
\caption{电源管理模块原理图（原理图）}
\label{fig:sch_power_module}
\end{figure}
电源模块实物与原理连接如图\ref{fig:photo_power}和图\ref{fig:sch_power_module}所示，后续电流预算如表\ref{tab:power_budget}。

\begin{table}[htbp]
\centering
\caption{电源分配与最大电流预算表}
\label{tab:power_budget}
\renewcommand{\arraystretch}{1.15}
\begin{tabular}{>{\raggedright\arraybackslash}p{2.1cm} >{\raggedright\arraybackslash}p{3.7cm} >{\centering\arraybackslash}p{2.2cm} >{\centering\arraybackslash}p{2.2cm} >{\raggedright\arraybackslash}p{2.6cm}}
\toprule
电源路 & 主要负载 & 典型电流(mA) & 最大预算(mA) & 备注 \\
\midrule
3.3\,V路 & STM32F103、DHT11温湿度链路、MQ-2采样侧、AS608逻辑接口、SD卡模块 & 120 & 260 & 写卡瞬态峰值纳入预算 \\
5\,V路 & ESP8266模块、PIR模块、部分外围接口 & 60 & 220 & 联网与外设供电 \\
12\,V路 & 门锁执行器、燃气阀驱动支路（继电器控制） & 0（待机）/120（动作） & 450 & 瞬态冲击按动作峰值预留 \\
可调路（0.9--29.7\,V） & 实验阶段外设兼容验证与临时负载 & 0--40 & 150 & 正式运行默认不带载 \\
\bottomrule
\end{tabular}
\end{table}

从表\ref{tab:power_budget}可以看到，系统在主工作状态下以3.3\,V路和5\,V路为主要功耗来源。其中，SD卡写入为短时突发负载（抓拍落盘/日志写入瞬间电流上升），因此在3.3\,V路预算中按峰值冗余预留。结合第6章表\ref{tab:perf_target_compare}中的整机功耗测试结果：待机功耗420\,mW、告警功耗980\,mW，可按$P=UI$折算为5\,V输入侧等效电流约84\,mA和196\,mA，均低于5\,V路220\,mA预算上限；告警工况下12\,V支路短时动作电流叠加后仍在整体预算范围内，说明电源分配与功耗测试数据一致，满足工程冗余要求。
\section{整机电路布局与接线汇总}
完成单模块电路说明后，还需要给出整机层面的接线关系。本文的实物搭建不是多块控制板分布式结构，而是以STM32F103RCT6核心板为中心，把传感、执行、通信和存储模块通过排线接入同一控制板。整机布局时优先把执行器和无线通信模块与模拟采样支路拉开距离，MQ-2采样线靠近地参考，SD卡和OV7670并口线尽量缩短，PIR和键盘放在便于触发的位置。这样布置可以减少继电器动作、ESP8266发射和长杜邦线对采样与抓拍链路的影响。

整机电路布局关系如图\ref{fig:whole_circuit_layout}所示。图中同时给出了主要信号引脚和供电分支，便于对照前文各模块原理图检查实物接线。
\begin{figure}[H]
\centering
\begin{tikzpicture}[
  node distance=9mm and 14mm,
  every node/.style={font=\zihao{5}, align=center},
  core/.style={rounded corners=3pt, draw=blue!70!black, fill=blue!8, line width=0.9pt, minimum width=34mm, minimum height=18mm},
  mod/.style={rounded corners=2pt, draw=gray!70!black, fill=gray!8, line width=0.75pt, minimum width=30mm, minimum height=10mm, inner ysep=3pt},
  pwr/.style={rounded corners=2pt, draw=orange!70!black, fill=orange!10, line width=0.75pt, minimum width=30mm, minimum height=10mm},
  sig/.style={-{Stealth[length=2.2mm]}, thick, draw=blue!65!black},
  pow/.style={-{Stealth[length=2.2mm]}, thick, draw=orange!70!black}
]
\node[core] (mcu) {STM32F103RCT6\\核心控制板};
\node[mod, above left=of mcu] (as608) {AS608指纹\\PA9/PA10, PA8};
\node[mod, left=of mcu] (pir) {PIR人体红外\\PC4/EXTI};
\node[mod, below left=of mcu] (mq2) {MQ-2燃气\\PB0/ADC};
\node[mod, above=of mcu] (dht) {DHT11温湿度\\PC5单总线};
\node[mod, above right=of mcu] (cam) {OV7670+FIFO\\PA0--PA7, PA11/PA12};
\node[mod, right=of mcu] (sd) {SD卡模块\\SPI2: PB12--PB15};
\node[mod, below right=of mcu] (esp) {ESP8266\\USART3: PB10/PB11};
\node[pwr, below=14mm of mcu] (power) {DCP3512电源\\3.3\,V/5\,V/12\,V};
\node[mod, below=16mm of power, minimum width=34mm, minimum height=15mm] (act) {执行器驱动\\继电器/蜂鸣器\\PC6/PC7/PC10};
\draw[sig] (as608) -- (mcu);
\draw[sig] (pir) -- (mcu);
\draw[sig] (mq2) -- (mcu);
\draw[sig] (dht) -- (mcu);
\draw[sig] (cam) -- (mcu);
\draw[sig] (sd) -- (mcu);
\draw[sig] (esp) -- (mcu);
\draw[sig] (mcu.south west) -- ++(-4mm,-6mm) |- (act.west);
\draw[pow] ([xshift=8mm]power.north) -- node[right]{3.3\,V} ([xshift=8mm]mcu.south);
\draw[pow] (power.west) -| node[pos=0.25,left]{5\,V} (mq2.south);
\draw[pow] (power.east) -| node[pos=0.25,right]{5\,V} (esp.south);
\draw[pow] ([xshift=8mm]power.south) -- node[right, midway]{12\,V执行支路} ([xshift=8mm]act.north);
\end{tikzpicture}
\caption{整机电路布局与主要接线关系}
\label{fig:whole_circuit_layout}
\end{figure}

为减少杜邦线直连带来的松动和错接风险，实物搭建阶段将各模块接口按功能分区汇总，形成整机原理图。该图把主控、HC-SR501、SD、ESP-01S、OLED、DHT11、MQ-2、蜂鸣器、OV7670+FIFO、4$\times$4矩阵键盘、继电器控制以及电源共地关系放在同一张图中，便于检查各模块接口和供电连接。整机原理图如图\ref{fig:system_overall_schematic}所示。

\begin{figure}[H]
\centering
\thesisfigfit{images/原理图/sch_system_overall.png}{0.98}{0.62}
\caption{整机原理图}
\label{fig:system_overall_schematic}
\end{figure}

\section{硬件联调与电气问题排查}
本节所述“硬件联调”不是只检查电路是否导通，而是围绕电源、地线、信号完整性、接口电平和外设连接可靠性来排查问题。调试阶段先完成供电、时钟、复位和共地检查，再借助万用表、示波器和逻辑分析仪确认UART、ADC、SPI、单总线和外部中断等接口的电气状态。与控制流程、任务调度和程序实现相关的内容放在第4章展开，本节只保留与硬件连接和电气稳定性直接相关的排查过程。  

结合本系统的实际联调过程，本文选取若干与实物模块直接相关的典型案例，按“问题现象$\rightarrow$排查过程$\rightarrow$解决方法$\rightarrow$验证结果”的顺序说明硬件问题如何定位到电源、接地、接口时序或器件参数。

\subsection*{典型联调闭环案例}
\kuan{1}案例A：继电器动作引起MQ-2采样抖动偏大  
\xiang{1}问题现象：继电器吸合/释放瞬间，MQ-2 ADC采样出现尖峰跳变，导致短时误触发风险上升。  
\xiang{2}排查过程：示波器分别观测3.3\,V数字电源与ADC参考点波形，并对继电器线圈两端关断瞬态进行抓波，发现尖峰与采样异常同步；同时排查接地回流路径，确认执行器回流与模拟采样地存在耦合。  
\xiang{3}解决方法：对MQ-2采样支路增加RC低通（与第3章信号调理一致），执行器与模拟采样进行电源分区与单点汇接；继电器线圈增加续流二极管并优化布线回流路径。  
\xiang{4}验证结果：复测中ADC尖峰明显收敛，阈值判定稳定性提升；系统级联动测试未出现由继电器动作引起的误报样本（与第6章误报率统计结论一致）。

\kuan{2}案例B：温湿度单总线采样偶发超时  
\xiang{1}问题现象：温湿度采样在部分时刻返回超时或校验失败，导致状态上报中断。  
\xiang{2}排查过程：使用逻辑分析仪抓取单总线时序，检查起始信号宽度、应答窗口与数据位脉宽；同时复核DATA线上拉电阻、模块供电和杜邦线长度，发现线缆较长时上升沿变缓。  
\xiang{3}解决方法：DATA线保持4.7\,k$\Omega$上拉，缩短传感器连线并保证模块与主控共地；必要时将信号线与地线并行布置，减小线缆寄生和外部干扰。  
\xiang{4}验证结果：连续采样稳定性恢复，单元测试中DHT11读取成功率达到100\%（如表\ref{tab:unit_results}）。

\kuan{3}案例C：PIR（HC-SR501）偶发误触发  
\xiang{1}问题现象：布防状态下出现“无人经过但触发”的偶发告警候选事件。  
\xiang{2}排查过程：检查PIR输出波形，发现存在窄脉冲毛刺；同时复核安装高度、朝向与电源噪声；对触发日志进行时间关联，定位部分样本与电源瞬态/强光干扰相关。  
\xiang{3}解决方法：保证PIR输入上拉与地参考稳定，调整模块安装朝向和灵敏度电位器，避开窗边强光、热源和气流直吹位置；供电端增加去耦电容，降低瞬态噪声对触发脚的影响。  
\xiang{4}验证结果：系统级测试入侵误报率维持在3.2\%（表\ref{tab:perf_target_compare}），满足$\le5.0\%$目标。

\kuan{4}案例D：ESP8266重连阶段串口丢帧/粘包  
\xiang{1}问题现象：网络波动后重连期间，云端命令偶发解析失败或执行回执丢失。  
\xiang{2}排查过程：使用逻辑分析仪观察USART3的TX/RX波形，检查波特率、空闲电平、线缆长度与共地状态；同时测量ESP8266发射瞬间的5\,V/3.3\,V电源波动，确认无线模块峰值电流会拉低局部供电。  
\xiang{3}解决方法：统一串口波特率并缩短TX/RX连线，保证ESP8266与主控可靠共地；在模块供电端增加储能和高频去耦电容，必要时单独走线供电，减少无线发射对串口电平的影响。  
\xiang{4}验证结果：远程控制成功率维持在高水平（见第6章联动测试），弱网样本主要表现为时延增加而非逻辑失效。

\kuan{5}案例E：执行器驱动“吸合但MCU复位”  
\xiang{1}问题现象：在蜂鸣器/继电器高频动作时，偶发出现MCU复位或显示短暂异常。  
\xiang{2}排查过程：示波器观测3.3\,V电源跌落与地弹噪声，发现执行器动作导致瞬态压降；检查去耦电容与供电走线，确认局部储能不足与回流干扰。  
\xiang{3}解决方法：在执行器电源端增加储能电容并优化去耦布局；继电器线圈保持续流二极管，强电流支路与MCU供电支路分开回流，降低执行器动作对主控电源的冲击。  
\xiang{4}验证结果：连续运行测试未再出现异常复位，72\,h稳定性测试通过（表\ref{tab:stability_record}）。

通过上述闭环案例，硬件联调不再只停留在现象修复，而是转向“机理定位+参数修正+复测验证”。这样处理后，问题复现路径更清楚，后续维护效率也更高。
硬件联调闭环流程如图\ref{fig:debug_closed_loop}所示。
\begin{figure}[H]
\centering
\begin{tikzpicture}[
  node distance=7mm,
  every node/.style={font=\zihao{-4}, align=center},
  flow/.style={
    rounded corners=2.5pt,
    draw=blue!55!black,
    fill=blue!8,
    line width=0.8pt,
    minimum width=4.2cm,
    minimum height=0.95cm,
    blur shadow={shadow blur steps=5, shadow blur extra rounding=1.5pt, shadow xshift=1pt, shadow yshift=-1pt, shadow opacity=15}
  },
  arrow/.style={-{Stealth[length=2.6mm]}, very thick, draw=blue!65!black}
]
\node[flow] (a) {问题发现};
\node[flow, below=of a] (b) {定位分析};
\node[flow, below=of b] (c) {方案修复};
\node[flow, below=of c] (d) {复测验证};
\node[flow, below=of d] (e) {稳定确认};
\node[font=\zihao{-5}, text=blue!55!black, right=1.6cm of e] (f) {反馈优化};

\draw[arrow] (a) -- (b);
\draw[arrow] (b) -- (c);
\draw[arrow] (c) -- (d);
\draw[arrow] (d) -- (e);
\draw[arrow] (e.east) to[out=0,in=0,looseness=1.15] (a.east);
\draw[arrow] (e.east) -- (f.west);
\end{tikzpicture}
\caption{硬件联调闭环图}
\label{fig:debug_closed_loop}
\end{figure}

\section{本章小结}
本章围绕主控最小系统、传感信号调理、执行驱动与隔离、电源管理、接口分配、整机电路布局和硬件联调完成硬件设计。通过参数计算、器件选型对比和抗干扰细节落实，方案在实时性、可靠性和可维护性之间保持了较好的工程平衡。硬件实现结果说明，单板集中架构可以支撑多外设并发接入，也能满足后续软件调度和联动控制需求，为第4章软件实现和第6章测试验证提供了硬件基础。

\chapter{系统软件设计}
\section{软件开发环境}
本项目的软件环境不是一次性确定，而是按“单模块可用—系统可并发—链路可闭环”的节奏逐步固定。前期先用Keil MDK和STM32外设库完成外设点亮、串口回显和中断验证；中期切换到FreeRTOS进行并发任务联调；后期再与ESP8266和OneNET打通远程链路。联调过程中主要依赖串口日志、命令回显和状态快照定位问题，以便将“外设问题”和“业务逻辑问题”分离处理。

\begin{table}[htbp]
\centering
\caption{软件开发环境与版本}
\renewcommand{\arraystretch}{1.15}
\begin{tabular}{>{\raggedright\arraybackslash}p{3.3cm} >{\raggedright\arraybackslash}p{2.6cm} >{\raggedright\arraybackslash}p{6.2cm}}
\toprule
软件/平台 & 版本 & 用途 \\
\midrule
Keil MDK & 5.x & 单芯片固件开发与调试 \\
STM32固件库 & 标准外设库 & 外设驱动接口支持 \\
FreeRTOS & V9.x & 全功能模式多任务调度 \\
OneNET MQTT & 3.1.1协议 & 云端消息中间件服务 \\
串口调试工具 & 任意串口助手 & JSON状态监控与命令调试 \\
\bottomrule
\end{tabular}
\end{table}
\section{软件架构}
软件部分按照“先可跑、再可并发、最后可维护”的思路组织。底层驱动统一硬件接口，上层业务不直接操作寄存器；服务层负责消息转发、状态同步和容错处理；业务层实现门禁、告警和应急逻辑；应用交互层负责云端和移动端对接\cite{freertos_book,freertos_sched}。这样拆分后，后续调整参数或替换模块时改动范围更小，也更符合毕设阶段逐步迭代的开发节奏。
软件总体流程如图\ref{fig:sw_overall_flow}所示。
\begin{figure}[H]
\centering
\thesisfigfit{images/fig4_0_sw_overall_flow.png}{0.92}{0.72}
\caption{软件总体流程图}
\label{fig:sw_overall_flow}
\end{figure}
软件架构与任务关系如图\ref{fig:sw_arch}所示。

\clearpage
\begin{figure}[p]
\centering
\vspace*{-0.2cm}
\thesisfigfit{images/fig4_1_sw_arch.png}{1.00}{0.95}
\caption{软件架构与任务关系图}
\label{fig:sw_arch}
\end{figure}
\clearpage

\begin{table}[htbp]
\centering
\caption{FreeRTOS任务划分与资源配置}
{\zihao{5}\renewcommand{\arraystretch}{1.2}
\begin{tabular}{>{\centering\arraybackslash}p{2.5cm} >{\centering\arraybackslash}p{1.4cm} >{\centering\arraybackslash}p{1.8cm} >{\centering\arraybackslash}p{1.8cm} >{\centering\arraybackslash}p{3.6cm}}
\toprule
任务名称 & 优先级值 & 堆栈（word） & 周期/触发 & 主要职责 \\
\midrule
\texttt{Task\_Lock} & 5 & 512 & 20\,ms周期 & 键盘扫描、门锁状态机与开锁控制 \\
\texttt{Task\_Net} & 4 & 1024 & 50\,ms周期 & ESP8266维护、MQTT解析、状态上报与重连 \\
\texttt{Task\_Sensor} & 3 & 256 & 100\,ms周期 & DHT11、PIR、MQ-2采样与告警联动 \\
\texttt{Task\_Log} & 3 & 1024 & 200\,ms周期 & 指纹轮询、录入状态机、抓拍轮询与日志同步 \\
\texttt{Task\_OLED} & 2 & 256 & 200\,ms周期 & 主界面刷新与状态提示 \\
\bottomrule
\end{tabular}
}
\end{table}
\noindent 堆栈分配依据采用“静态估算+高水位校核”方法：先按任务局部变量、调用深度和通信缓冲区估算初值，再通过\texttt{uxTaskGetStackHighWaterMark()}在线观测，保证最小剩余栈空间大于20\%。

\subsection{任务堆栈大小的估算依据}
以Cortex-M3为例，任务堆栈需求可近似拆分为“上下文保存+函数调用深度+局部变量与缓冲区+安全裕量”。工程上采用下式给出初始估算：
\begin{equation}
S_{\mathrm{word}} \ge \frac{S_{\mathrm{ctx}} + S_{\mathrm{call}} + S_{\mathrm{local}} + S_{\mathrm{buf}}}{4}\cdot (1+\eta)
\label{eq:stack_budget}
\end{equation}
其中$S_{\mathrm{ctx}}$为任务切换上下文开销字节数，$S_{\mathrm{call}}$为最大调用深度对应的栈帧开销，$S_{\mathrm{local}}$与$S_{\mathrm{buf}}$分别为局部变量和通信缓冲区开销，$\eta$为裕量系数（工程建议$\eta\ge 0.2$）。初值配置后通过高水位统计进行校核并迭代收敛，避免“堆栈过小导致HardFault”与“过大挤占RAM”两类问题。

\noindent 任务间通信与同步机制采用“队列+事件组+互斥锁”的组合机制，具体安排如下：  
\kuan{1}队列（Queue）：用于跨任务传递传感采样帧和告警事件，典型队列深度为8，消息结构体固定长度，避免堆碎片；  
\kuan{2}事件组（EventGroup）：用于布防状态变化、网络在线状态、抓拍触发等位标志同步；  
\kuan{3}二值信号量（Binary Semaphore）：用于中断到任务的快速唤醒，如PIR触发后置位告警事件并由传感/联动轮询路径处理；  
\kuan{4}互斥锁（Mutex）：用于共享资源保护，如串口发送缓存和日志文件写入，防止并发写冲突。

硬实时任务响应保障采用三项措施：中断仅置位事件不做耗时处理；\texttt{Task\_Lock}保持最高优先级并禁止阻塞式长延时；网络与日志任务采用分片执行，避免抢占关键执行窗口。

\noindent 共享资源保护和ISR安全接口按“互斥锁优先、临界区最小化”处理：  
\kuan{1}互斥锁：跨任务访问共享资源优先使用Mutex，避免长时间关中断；  
\kuan{2}临界区：仅在更新少量共享标志或环形缓冲指针时使用\texttt{taskENTER\_CRITICAL()}与\texttt{taskEXIT\_CRITICAL()}，并严格控制临界区长度；
\kuan{3}ISR到任务：中断中使用\texttt{xSemaphoreGiveFromISR}/\texttt{xQueueSendFromISR}等ISR安全API，并在需要时触发上下文切换，保证硬实时链路及时唤醒。  
该策略与本论文“中断只做快判定、复杂处理下沉任务”的原则一致，可在保证实时性的同时降低死锁与优先级反转风险。裸机轮询路径不再单独作为一节展开：当裸机模式启用时，系统不启动调度器，而由主循环分段执行键盘、指纹、安防和串口控制等逻辑，PIR等快速事件仍由中断置位。该路径主要服务于硬件联调和故障定位；进入完整联网与并发运行后，再切换至FreeRTOS模式。
\noindent 底层驱动开发阶段，本文将目标限定为接口统一、时序可控、异常可恢复和模块可复用。为减少业务层和硬件细节耦合，驱动层统一采用“初始化-数据收发-状态查询-异常处理”接口范式，覆盖GPIO、UART、I2C/单总线、定时器和外部中断等模块。按这一范式组织后，即使后续更换器件或升级模块，上层控制逻辑通常也不需要大改，兼容性更稳定。
其中SD驱动实施为SPI2专用通道（PB13/PB14/PB15+PB12片选），并将初始化过程显式分为“上电稳定等待$\rightarrow$发送80字节空时钟$\rightarrow$CMD/ACMD握手$\rightarrow$FatFs挂载”四步。针对低质量卡或线缆寄生影响，驱动中保留较宽松的命令等待和ACMD41重试参数（如200\,ms命令等待窗口、600次轮询上限），优先保证挂载成功率；仅在挂载成功后开放抓拍文件写入与容量查询接口，避免业务层在未就绪状态下访问存储造成异常。
\section{核心控制逻辑}
核心控制逻辑按门禁、安防、环境和场景四类业务拆分实现，并通过有限状态机（FSM）和规则引擎完成事件判定与动作调度。策略上坚持“本地闭环优先、远程协同增强”，确保网络不稳定时基础安全功能仍可独立运行。
门禁与安防核心状态机如图\ref{fig:state_machine}所示。

\begin{figure}[H]
\centering
\thesisfigfit{images/fig4_2_state_machine.png}{0.88}{0.72}
\caption{门禁与安防核心状态机}
\label{fig:state_machine}
\end{figure}

\subsection{门禁验证逻辑}
在门禁实现中，本文将状态机定义为\texttt{IDLE}、\texttt{ACQUIRE}、\texttt{MATCH}、\texttt{OPEN}、\texttt{LOCKOUT}五个状态。状态转移条件如下：  
\kuan{1}\texttt{IDLE$\rightarrow$ACQUIRE}：检测到触摸/按键触发；  
\kuan{2}\texttt{ACQUIRE$\rightarrow$MATCH}：采集成功且质量分大于阈值；  
\kuan{3}\texttt{MATCH$\rightarrow$OPEN}：模板匹配成功；  
\kuan{4}\texttt{MATCH$\rightarrow$LOCKOUT}：连续失败次数$\geq N_f$（本系统取$N_f=3$）；  
\kuan{5}\texttt{LOCKOUT$\rightarrow$IDLE}：锁定计时到期或管理员复位。  
\noindent 门禁验证流程如图\ref{fig:door_auth_flow}所示。
\begin{figure}[H]
\centering
\begin{tikzpicture}[
  node distance=8mm and 14mm
]
\node[flowblock] (start) {触发输入\\(按键/指纹)};
\node[flowblock, below=of start] (acq) {采集凭据\\(\texttt{ACQUIRE})};
\node[flowdecision, below=of acq] (qok) {质量分\\达标?};
\node[flowblock, below=of qok] (match) {模板匹配\\(\texttt{MATCH})};
\node[flowdecision, below=of match] (mok) {匹配\\成功?};
\node[flowblock, below=of mok] (open) {开锁执行\\(\texttt{OPEN})};
\node[flowblock, below=of open] (idle) {返回待机\\(\texttt{IDLE})};
\node[flowblock, right=26mm of mok] (fail) {失败计数+1};
\node[flowdecision, below=of fail] (lock) {$fail\_cnt \ge N_f$?};
\node[flowblock, below=of lock] (lockout) {进入锁定\\(\texttt{LOCKOUT})};

% 左侧汇流竖线 x 坐标：所有“否”回绕走外侧，避免穿过主链节点
\coordinate (busX) at ($(qok.west)+(-2.65,0)$);

\draw[flowline] (start) -- (acq);
\draw[flowline] (acq) -- (qok);
\draw[flowline] (qok) -- node[flowlabel, right]{是} (match);
\draw[flowline] (match) -- (mok);
\draw[flowline] (mok) -- node[flowlabel, right]{是} (open);
\draw[flowline] (open) -- (idle);
\draw[flowline] (mok.east) -- node[flowlabel, above]{否} (fail.west);
\draw[flowline] (fail) -- (lock);
\draw[flowline] (lock) -- node[flowlabel, right]{是} (lockout);
\draw[flowline] (qok.west) -- node[flowlabel, above, pos=0.18]{否} ($(qok.west)+(-0.5,0)$) -- (busX |- qok.west) -- (busX |- idle.west) -- (idle.west);
\draw[flowline] (lock.west) -- node[flowlabel, above, pos=0.15]{否} ($(lock.west)+(-0.5,0)$) -- (busX |- lock.west) -- (busX |- acq.west) -- (acq.west);
\draw[flowline] (lockout.west) -- ++(-1.35,0) coordinate (loH)
  -- (loH |- idle.south) -- (idle.south east);
\node[flowlabel, left=1mm of loH, text width=16mm, align=center] {计时到期\\管理员复位};
\end{tikzpicture}
\caption{门禁验证逻辑流程图}
\label{fig:door_auth_flow}
\end{figure}
边界异常处理在固件中实现为：串口超时最多重试2次、模板库访问异常回退到\texttt{IDLE}并记录错误日志、门锁执行反馈超时时触发二次驱动并上报故障。
\subsection{安防告警逻辑}
在安防链路中，本文将状态机定义为\texttt{DISARMED}、\texttt{ARMED}、\texttt{TRIGGERED}、\texttt{ALARMING}、\texttt{RECOVERY}。其中：  
\kuan{1}\texttt{ARMED$\rightarrow$TRIGGERED}：PIR触发并通过最小持续时间门限；  
\kuan{2}\texttt{TRIGGERED$\rightarrow$ALARMING}：抓拍轮询路径返回本地落盘结果或超时兜底；  
\kuan{3}\texttt{ALARMING$\rightarrow$RECOVERY}：用户确认解除或风险消失持续$T_{clr}$。  
\noindent 安防告警流程如图\ref{fig:alarm_logic_flow}所示。
\begin{figure}[H]
\centering
\begin{tikzpicture}[
  node distance=7mm and 11mm
]
\node[flowdecision] (arm) {布防状态?};
\node[flowblock, below=of arm] (pir) {PIR触发并通过\\最小持续门限};
\node[flowblock, below=of pir] (trg) {进入\texttt{TRIGGERED}\\触发抓拍/复核};
\node[flowblock, below=of trg] (alm) {本地告警执行\\蜂鸣+日志+上报};
\node[flowblock, below=of alm] (rec) {解除条件满足\\进入\texttt{RECOVERY}};
\node[flowblock, below=of rec] (back) {返回\texttt{ARMED}/\texttt{DISARMED}};
\node[flowblock, right=28mm of arm] (ignore) {忽略触发\\保持\texttt{DISARMED}};

\draw[flowline] (arm) -- node[flowlabel, right]{是} (pir);
\draw[flowline] (pir) -- (trg);
\draw[flowline] (trg) -- (alm);
\draw[flowline] (alm) -- (rec);
\draw[flowline] (rec) -- (back);
\draw[flowline] (arm.east) -- node[flowlabel, above]{否} (ignore.west);
\end{tikzpicture}
\caption{安防告警逻辑流程图}
\label{fig:alarm_logic_flow}
\end{figure}
异常边界策略为：抓拍失败不阻塞本地告警；云端不可达时保留本地蜂鸣、SD落盘和日志；执行器反馈异常时进入降级模式并限制高风险动作。
\subsection{环境监测与联动}
系统周期采集DHT11温湿度和MQ-2气体数据，采用“链路稳态保障+数字滤波+趋势判定+阈值触发”联合策略。该策略的设计目标是：在家庭环境存在温湿波动、油烟干扰和电磁噪声的条件下，仍能稳定输出可用于联动控制的测量值，并将误触发控制在第6章给出的指标范围内。

\subsubsection*{DHT11链路稳定性设计}
DHT11温湿度数据主要用于环境状态展示与联动阈值辅助判定。为降低误读对联动逻辑的影响，系统采用“周期采样+异常重试+超时回退”的处理方式。该链路在测试阶段保持了稳定采样和状态更新能力，能够满足第6章对状态展示与联动稳定性的要求。

\subsubsection*{MQ-2标定、ADC到浓度转换模型与温漂补偿}
MQ-2输出为与气体浓度相关的模拟量，本系统以ADC采样值$A_k$作为原始观测量。结合本科工程实现难度与家庭场景需求，本文采用“两阶段标定”策略：  
\kuan{1}阈值标定：在安全环境下采集稳定基线$A_{\mathrm{base}}$，在模拟超限场景下采集$A_{\mathrm{alarm}}$，据此确定阈值$T_{gas}$与变化率阈值$\Delta T$；  
\kuan{2}可选浓度换算：在具备标准气源或半定量参考仪表条件时，建立ADC与浓度的工程映射模型，用于更直观的浓度展示与跨环境对比。

在可选浓度换算中，为兼顾可实现性与拟合效果，采用分段线性模型：
\begin{equation}
C=
\begin{cases}
k_1(A-A_{\mathrm{base}}), & A_{\mathrm{base}}\le A<A_1\\
k_2(A-A_1)+C_1, & A_1\le A<A_2\\
k_3(A-A_2)+C_2, & A\ge A_2
\end{cases}
\label{eq:mq2_piecewise}
\end{equation}
其中$k_1,k_2,k_3$分别为三段斜率系数（单位ppm/ADC），$A_1,A_2$为分段拐点；$C_1=k_1(A_1-A_{\mathrm{base}})$、$C_2=k_2(A_2-A_1)+C_1$用于保证分段连接连续。温漂补偿采用“温湿修正+基线跟踪”组合：一方面对阈值作小幅修正$T_{gas}(T,H)$；另一方面在无告警事件且变化率较小的窗口内对基线缓慢更新：
\begin{equation}
A_{\mathrm{base}}(k)=\lambda A_{\mathrm{base}}(k-1)+(1-\lambda)\bar{A}_k,\quad 0<\lambda<1
\label{eq:mq2_baseline_track}
\end{equation}
该策略用于抑制预热漂移与长期老化导致的阈值失配，提升长期稳定性（与第6章72\,h稳定性测试结论一致）。

\subsubsection*{中值滤波/滑动平均滤波实现逻辑与降噪机理}
为抑制继电器动作、电源纹波和瞬态干扰引起的尖峰噪声，MQ-2通道采用“中值滤波+滑动平均”级联结构：先用中值滤波剔除脉冲异常，再用滑动平均降低随机噪声。中值滤波窗口为$2r+1$：
\begin{equation}
\tilde{A}_k=\mathrm{median}\{A_{k-r},\ldots,A_{k+r}\}
\label{eq:median_filter}
\end{equation}
滑动平均为
\begin{equation}
\bar{A}_k=\frac{1}{M}\sum_{i=0}^{M-1}\tilde{A}_{k-i}
\label{eq:mean_filter}
\end{equation}
再计算变化率$\Delta A_k=\bar{A}_k-\bar{A}_{k-1}$。当$\bar{A}_k>T_{gas}$且$\Delta A_k>\Delta T$时判定为高风险事件。滤波前后对比数据如表\ref{tab:mq2_filter_compare}所示。

\begin{table}[htbp]
\centering
\caption{MQ-2滤波前后ADC数据对比（节选）}
\label{tab:mq2_filter_compare}
\renewcommand{\arraystretch}{1.15}
\begin{tabular}{>{\centering\arraybackslash}p{1.4cm} >{\centering\arraybackslash}p{2.3cm} >{\centering\arraybackslash}p{2.3cm} >{\centering\arraybackslash}p{2.3cm} >{\centering\arraybackslash}p{2.3cm}}
\toprule
样本序号 & 原始ADC $A_k$ & 中值输出$\tilde{A}_k$ & 均值输出$\bar{A}_k$ & 说明 \\
\midrule
1 & 214 & 214 & 214.0 & 正常基线 \\
2 & 219 & 219 & 216.5 & 轻微波动 \\
3 & 286 & 221 & 218.0 & 尖峰被抑制 \\
4 & 223 & 223 & 219.3 & 恢复稳定 \\
5 & 217 & 217 & 218.8 & 正常基线 \\
6 & 312 & 224 & 220.8 & 干扰尖峰被抑制 \\
\bottomrule
\end{tabular}
\end{table}

由表\ref{tab:mq2_filter_compare}可见，级联滤波可有效抑制继电器动作等瞬态尖峰，减少阈值误触发。实测标准差由约31.4\,ADC降至约9.7\,ADC（下降约69.1\%）。该实现逻辑与第6章“MQ-2气体检测成功率94.0\%”的单元测试结果一致，滤波与基线补偿用于降低误触发并提升阈值判定稳定性。
\subsection{场景联动逻辑}
本文在场景联动中实现了“居家”“离家布防”“夜间安防”等预设模式，并支持移动端参数配置。联动执行采用“条件触发+动作链”机制，在不增加硬件复杂度的前提下实现设备自动协同。
\noindent 场景联动执行流程如图\ref{fig:scene_linkage_flow}所示。
\begin{figure}[H]
\centering
\begin{tikzpicture}[
  node distance=7mm and 10mm
]
\node[flowblock] (mode) {读取场景模式\\(居家/离家/夜间)};
\node[flowblock, below=of mode] (sample) {采集门禁、PIR、MQ-2\\与环境状态};
\node[flowdecision, below=of sample] (judge) {规则匹配\\(阈值+时间窗)};
\node[flowblock, below=of judge] (act) {执行动作链\\(门锁/蜂鸣/阀门/抓拍)};
\node[flowblock, below=of act] (report) {状态回写与日志记录\\(本地优先，云端协同)};
\node[flowblock, below=of report] (loop) {进入下一周期轮询};

\draw[flowline] (mode) -- (sample);
\draw[flowline] (sample) -- (judge);
\draw[flowline] (judge) -- (act);
\draw[flowline] (act) -- (report);
\draw[flowline] (report) -- (loop);
\end{tikzpicture}
\caption{场景联动逻辑流程图}
\label{fig:scene_linkage_flow}
\end{figure}

\subsection{HC-SR501距离校准、误触发抑制与联合判定}
HC-SR501作为入侵触发源，其误触发会直接影响系统误报率指标。本文采用“距离-视场校准+软件防抖+多条件联合判定”三步抑制误报：  
\kuan{1}距离校准：在走廊/门厅场景设置1\,m、3\,m、5\,m等点位，分别统计触发率，确定灵敏度电位器的工程配置区间，并记录最佳安装高度与朝向；  
\kuan{2}软件防抖：对PIR输入采用最小持续时间门限$T_{\min}$与锁定窗口$T_{\mathrm{lock}}$。当高电平持续时间小于$T_{\min}$判为毛刺丢弃；触发成立后在$T_{\mathrm{lock}}$内忽略重复触发，避免抖动造成重复上报；  
\kuan{3}联合判定：在布防状态下，PIR触发仅作为“候选事件”，需满足附加条件（例如二次触发、图像抓拍返回有效帧、或在短时间窗内的复核信号）才升级为入侵告警。  
该策略与第6章入侵检测误报率3.2\%结果对应，即通过“PIR触发+复核”机制将短时干扰从告警链路中剔除。

\subsection{多传感器融合的安防事件判定算法}
为进一步降低误报并提高鲁棒性，本文采用证据加权方式进行事件融合判定。设在时间窗$W$内，PIR触发证据为$e_{pir}\in\{0,1\}$，图像有效帧证据为$e_{cam}\in\{0,1\}$，门禁异常证据为$e_{acc}\in\{0,1\}$，环境异常证据为$e_{env}\in\{0,1\}$。定义融合评分：
\begin{equation}
S = w_{pir}e_{pir}+w_{cam}e_{cam}+w_{acc}e_{acc}+w_{env}e_{env},\qquad \sum w_i=1
\label{eq:fusion_score}
\end{equation}
当$S\ge \theta$时判定为高可信事件并进入告警状态；当$S<\theta$时仅记录日志并保持监测。权重分配遵循“可置信证据优先”原则：本地抓拍落盘结果与门禁异常的权重高于单次PIR触发；环境异常与入侵告警分属不同风险域，通常不直接互相触发，仅用于提升整体风险等级。工程实现中，该融合策略对应事件组与队列：PIR中断置位$e_{pir}$并唤醒安防任务，抓拍任务回传本地落盘结果$e_{cam}$，门禁任务更新$e_{acc}$，环境任务更新$e_{env}$，最终由安防任务在窗口$W$内完成融合判定。  
融合效果的量化证明以系统级实测指标为准：第6章给出的入侵检测误报率3.2\%即为融合与复核策略综合作用下的结果，满足$\le 5.0\%$目标。
系统功能流程如图\ref{fig:function_flow}所示。

\begin{figure}[htbp]
\centering
\thesisfigfit{images/fig4_4_function_flow.png}{0.94}{0.82}
\caption{系统功能流程}
\label{fig:function_flow}
\end{figure}
\section{串口与MQTT数据协议设计}
通信协议按“本地可调试、远程可协同”两条链路并行设计。本地串口链路优先保证可读和易排障，用于输出状态、查询日志和回放异常；云端链路通过MQTT承载告警与控制命令，重点保证弱网下的重连与去重。这样做的原因是：开发阶段先靠串口把功能跑通，再把同一份状态映射到云端，联调成本更低。
设备与云端通信时序如图\ref{fig:comm_timing}所示。

\begin{figure}[H]
\centering
\thesisfig{images/fig4_3_comm_timing.png}{0.90}
\caption{设备与云端通信时序}
\label{fig:comm_timing}
\end{figure}
\subsection{帧格式}
通信帧由帧头、命令字、长度、数据域与校验字段构成。帧头用于同步识别，长度字段用于动态解析数据域，校验字段用于检错，提高数据传输可靠性。

\begin{table}[htbp]
\centering
\caption{本地串口JSON字段定义}
\begin{tabular}{>{\centering\arraybackslash}p{2.4cm} >{\centering\arraybackslash}p{2.2cm} >{\centering\arraybackslash}p{2.2cm} >{\centering\arraybackslash}p{4.8cm}}
\toprule
字段名 & 类型 & 典型值 & 说明 \\
\midrule
temp & float & 25.0 & 温度（DHT11） \\
humi & float & 60.0 & 湿度（DHT11） \\
gas & int & 1200 & MQ-2 ADC值 \\
door & string & closed/open & 门锁状态 \\
arm & int & 0/1 & 撤防/布防状态 \\
sd.status & string & \makecell{mounted/\\unmounted} & SD挂载状态 \\
\bottomrule
\end{tabular}
\end{table}
\subsection{指令集}
命令接口分为串口日志命令与云端控制命令。串口侧采用统一日志查询命令格式查询历史记录；云端侧通过MQTT下发开锁、布防与撤防等控制指令，设备执行后回传状态。

\begin{table}[htbp]
\centering
\caption{通信命令接口}
{\zihao{5}\renewcommand{\arraystretch}{1.2}
\begin{tabular}{>{\centering\arraybackslash}p{2.7cm} >{\centering\arraybackslash}p{2.0cm} >{\centering\arraybackslash}p{2.5cm} >{\centering\arraybackslash}p{4.6cm}}
\toprule
接口名称 & 典型指令 & 方向 & 功能说明 \\
\midrule
\makecell{串口状态上报} & JSON行文本 & 设备$\rightarrow$PC & 周期输出设备状态 \\
\makecell{日志查询命令} & \makecell{CMD:LOG=\\LAST,10} & PC$\rightarrow$设备 & 查询历史日志 \\
\makecell{MQTT属性上报} & OneNET数据点 & 设备$\rightarrow$云端 & 上传温湿度/状态 \\
\makecell{MQTT控制命令} & 开锁/布防指令 & 云端$\rightarrow$设备 & 远程控制执行 \\
\bottomrule
\end{tabular}
}
\end{table}
\subsection{MQTT主题与消息设计}
为避免业务耦合，系统按OneNET物模型主链路组织MQTT消息：设备侧通过属性上报通道周期发布状态，云端通过属性下发通道发送控制命令。当前验收阶段聚焦门锁、布防、传感和本地抓拍落盘状态；图片链路保持“本地SD优先”，不将图片上传作为通过条件。  
\kuan{1}主链路：使用OneNET物模型\texttt{thing/property/post}与\texttt{thing/property/set}完成状态上报与控制下发；  
\kuan{2}事件链路：告警事件通过统一事件上报通道携带事件类型与时间戳；  
\kuan{3}图片增强链路（后续）：若启用方案二，可在非\texttt{\$sys}前缀的自定义主题上传图片分包，并由服务端重组后提供HTTP访问；  
\kuan{4}命令回执：控制执行后返回标准回执字段，确保“下发—执行—确认”闭环可追踪。  
控制命令统一包含\texttt{cmdId}、\texttt{ts}、\texttt{nonce}、\texttt{payload}、\texttt{crc16}等字段，设备端先做校验与去重，再进入业务执行流程。
\subsection{校验与容错}
协议采用CRC校验、应答超时重发、幂等去重和异常回退机制处理误码与丢包问题。串口侧采用CRC16-Modbus校验，MQTT侧采用\texttt{cmdId+ts}实现幂等控制。系统对非法帧、重复帧与超时帧设置容错策略，确保干扰环境下链路稳定。

\begin{table}[htbp]
\centering
\caption{通信容错策略}
\begin{tabular}{>{\centering\arraybackslash}p{3.2cm} >{\centering\arraybackslash}p{3.0cm} >{\centering\arraybackslash}p{4.8cm}}
\toprule
异常类型 & 判定条件 & 处理策略 \\
\midrule
校验失败 & Check不一致 & 丢弃帧并计数 \\
应答超时 & 超过T\_ack未收到ACK & 最多重发3次 \\
重复帧 & 序号重复（可选） & 去重处理，避免重复执行 \\
非法命令 & Cmd不在指令表 & 返回ERR并记录日志 \\
\bottomrule
\end{tabular}
\end{table}
\section{关键时序设计}
在STM32这类资源受限平台上，系统稳定性很大程度取决于时序安排，而不是某一个单点算法有多复杂。结合实测结果，本文把任务分成两类：一类是必须尽快完成的安全动作链路，包括门禁、入侵和燃气应急；另一类是允许短时延后的服务链路，包括状态上报、日志整理和远程同步。并发冲突出现时，系统优先保证前一类任务，后一类任务可以排队等待，从而避免非关键任务挤占关键执行窗口。
任务调度时序关系如图\ref{fig:task_gantt}所示。

\begin{figure}[H]
\centering
\thesisfig{images/fig4_5_task_gantt.png}{0.88}
\caption{任务调度甘特图（含优先级抢占与阻塞点）}
\label{fig:task_gantt}
\end{figure}

\kuan{1}任务周期与优先级协同：FreeRTOS模式下，各任务按“安全优先、通信次之、展示最后”的原则设置优先级。\texttt{Task\_Lock}使用20\,ms短周期保证键盘和门锁状态机响应，\texttt{Task\_Net}采用50\,ms周期维护网络链路，\texttt{Task\_Sensor}和\texttt{Task\_Log}同级轮转以兼顾传感采样、指纹处理和抓拍轮询，\texttt{Task\_OLED}低频刷新。这样配置的目的，是避免高频但非关键任务长期占用调度窗口，拖慢关键告警动作。测试中可见，多事件并发时低优先级任务会有轻微排队，但门禁与应急控制链路仍在目标时延内完成，说明当前优先级配置和任务划分是有效的。

\kuan{2}中断与轮询边界：系统采用“中断只做快判定、复杂处理下沉任务”的工程策略。PIR触发等突发事件在中断上下文中仅完成状态置位和时间戳记录，不执行耗时逻辑；后续联动判断与动作控制由任务上下文完成。该设计可减少中断占用时间、降低嵌套冲突概率，并提升系统整体可维护性。裸机模式下同样遵循这一原则，通过主循环分段执行替代阻塞式串行处理，确保关键输入源不被长期饥饿。

\kuan{3}缓冲区与数据一致性：针对串口、MQTT与日志模块并行访问共享数据的问题，系统采用“快照上报+最小临界区”方式降低数据竞争。状态上报前先复制关键字段形成只读快照，再交由通信模块编码与发送；日志写入采用固定格式与长度约束，避免动态拼接造成堆栈压力。对于网络重发场景，引入消息序号或时间标签用于去重判定，防止重复控制命令导致执行器误动作。

\kuan{4}故障降级与恢复：考虑家庭网络波动与传感器偶发异常，系统实现了面向场景的降级机制。云端链路不可用时自动降级为本地自治模式；单个传感器异常时维持其余链路继续工作并记录日志；执行器反馈超时或状态不一致时触发重试与告警提示。该机制将“局部异常”限制在局部模块内，避免扩散为“全局失效”。从测试结果看，降级与恢复路径均可被稳定触发，验证了系统在工程运行中的韧性设计。
异常处理与恢复流程如图\ref{fig:fault_recovery}所示。

\begin{figure}[htbp]
\centering
\thesisfigfit{images/fig4_6_fault_recovery.png}{0.90}{0.78}
\caption{异常处理与恢复流程图}
\label{fig:fault_recovery}
\end{figure}

\subsubsection*{（1）传感器离线检测与恢复}
为避免“传感器异常$\rightarrow$错误联动”，系统对关键传感器建立离线检测与恢复策略：  
\kuan{1}DHT11温湿度链路：若连续$N_{to}$次采样超时或校验失败，则置离线标志并冻结温湿度显示为“无效”；进入离线后采用指数退避策略降低重试频率，避免占用CPU；当连续$N_{ok}$次采样成功则自动恢复在线并清除离线标志。  
\kuan{2}MQ-2：若采样长期恒定（变化率低且持续时间超过门限）或超出ADC合理范围，则判为异常；异常期间仅记录日志且禁止触发应急关阀动作，防止误联动；恢复条件为采样恢复到合理区间并通过滤波稳定窗口。  
\kuan{3}PIR：若在非布防状态出现高频触发，则判为干扰；进入抑制窗口$T_{\mathrm{sup}}$并记录环境干扰日志，避免无意义上报。

\subsubsection*{（2）ESP8266断网自动重连与命令幂等}
网络链路采用“状态机+分级重试”策略维护连接：  
\kuan{1}链路检测：周期发送心跳或PING，连续失败计数达到阈值后标记离线；  
\kuan{2}重连流程：执行\texttt{WiFi Join}$\rightarrow$\texttt{MQTT Connect}$\rightarrow$\texttt{Subscribe}$\rightarrow$\texttt{Sync State}；  
\kuan{3}重试策略：采用指数退避（如1\,s、2\,s、4\,s…上限不超过60\,s）降低弱网拥塞；  
\kuan{4}幂等去重：对下行控制命令以\texttt{cmdId+ts}为键做去重，避免重连/重发造成执行器重复动作；对执行结果上报采用回执确认，保证“下发-执行-确认”闭环。

\subsubsection*{（3）看门狗喂狗与异常复位恢复}
为提高长时间运行稳定性，系统引入独立看门狗（IWDG）作为最后保护：  
\kuan{1}喂狗策略：由\texttt{Task\_Log}统一执行总喂狗，喂狗前检查关键任务心跳（各任务周期性更新计数器或时间戳）；若任一关键任务超过阈值未更新，则禁止喂狗，触发看门狗复位，实现“故障自愈”；  
\kuan{2}复位恢复：启动时读取复位源（上电复位/看门狗复位），若为看门狗复位则进入安全模式：执行器置为安全态（关阀、门锁保持默认态、蜂鸣器关闭），并上报“异常复位”日志；随后逐步恢复网络与业务任务，避免复位后立即进入同一故障循环；  
\kuan{3}异常捕获：对HardFault等不可恢复异常，记录最小化故障上下文（关键寄存器/错误码）到日志区，为后续定位提供依据。
\section{网络与安全机制}
根据远程互联设计要求，网络通信方案需满足低带宽开销、断线可恢复、跨平台兼容与安全可加固等条件。考虑家庭网络波动与公网传输风险，系统采用“应用层轻量协议+传输层加密+本地敏感数据保护”的分层安全策略，在保证实时通信的同时降低数据泄露与非法控制风险。
\kuan{1}WiFi与MQTT链路：根据设备联网需求，通信链路应支持低功耗终端接入、主题化消息分发与弱网重连机制。综合报文开销、实现复杂度与生态成熟度，本设计选用ESP8266作为WiFi接入模块，并采用MQTT发布/订阅模型完成状态上报与控制下发，满足家庭场景下低带宽、可扩展的远程通信要求\cite{mqtt_spec,mqtt_v5,esp8266_at}。  
\kuan{2}AES-128加密存储：根据本地数据安全要求，系统采用AES-128对敏感信息进行分块加密存储\cite{aes_fips197}。实现流程为：  
\xiang{1}密钥装载：设备启动时从受保护区域加载主密钥并派生会话密钥；  
\xiang{2}数据封装：按16字节分组并填充，附加版本号与时间戳；  
\xiang{3}分组加密：对配置参数、指纹索引和关键日志执行AES-128-CBC加密；  
\xiang{4}完整性校验：存储密文同时保存HMAC/CRC校验值；  
\xiang{5}解密恢复：读取时先校验完整性，再解密并回滚异常记录。  
密钥轮换采用“版本号+双密钥窗口”策略，保证升级阶段兼容旧数据。
\kuan{3}HTTPS/TLS与MQTT鉴权：根据远程协同控制安全要求，传输链路需具备机密性、完整性与身份认证能力。App到云平台的接口采用HTTPS/TLS，保障移动端访问过程中的数据安全；设备端通过Token鉴权接入OneNET MQTT TCP链路，满足当前状态上报与远程控制需求。系统在ESP8266侧预留MQTT TLS初始化与证书校验路径，支持后续升级到服务器证书校验、随机数握手与会话复用\cite{rfc5246,tls_rfc}。在资源受限场景下，采用“当前链路先稳定、TLS链路可扩展”的策略平衡安全与时延。

\section{本章小结}
本章从分层软件架构、FreeRTOS任务资源配置、状态机控制逻辑、多传感器融合算法、通信协议和安全机制六个方面说明了软件实现路径。系统通过“中断快响应+任务化执行+分级容错”保证关键路径实时性，同时兼顾联网能力和后期维护；通过AES-128本地加密、App侧HTTPS/TLS和设备端Token鉴权，提高存储与通信链路的安全性。上述设计和第6章量化测试指标相互对应，形成了可复现、可验证的软件闭环。

\chapter{移动端交互设计与实现}
\section{移动端总体设计}
本系统使用\texttt{uni-app}开发移动端页面，将其作为主要交互入口，并通过OneNET云平台完成远程监测和控制。移动端不直接替代本地安防逻辑，而是承担状态可视化、远程命令下发和异常结果查看三项职责。当前收尾阶段先完成状态显示、远程开锁、布撤防和抓拍指令入口，图片在线展示留到HTTP图床或MQTT分包链路稳定后再纳入验收。

移动端交互采用“页面展示层+云端通信层+设备响应层”结构：页面层负责首页总览、控制按钮和抓拍入口；通信层通过OneNET REST接口查询属性并下发控制命令；设备侧按物模型解析命令，执行后再回传状态\cite{mqtt_iot_review}。这样划分后，移动端页面改动不会影响设备端本地闭环，设备离线时也不会阻断门禁、燃气和入侵告警等关键安全链路。

界面设计围绕“首页总览—云端抓拍”展开。首页集中展示温度、湿度、燃气、门锁、布防、报警和LED状态，并提供刷新、远程开锁、布防模式切换、报警阈值保存和摄像头查看入口；云端抓拍页面提供物模型刷新、手动抓拍和图片显示区域。图\ref{fig:mobile_home}和图\ref{fig:mobile_camera}分别给出两个主要页面。

如图\ref{fig:mobile_home}所示，首页顶部显示OneNET产品与设备信息，并给出平台返回值和更新时间，用于判断REST查询是否成功。中部用卡片展示温度、湿度、燃气、报警、门锁和布防状态；下方提供LED灯控、远程开锁、布防模式切换以及燃气和温湿度阈值保存功能。页面底部保留摄像头抓拍入口，说明当前验收阶段图片以本地SD保存为主，App页面只预留查看位置。

\begin{figure}[H]
\centering
\thesisfig{首页.jpg}{0.42}
\caption{移动端首页状态与控制界面}
\label{fig:mobile_home}
\end{figure}

如图\ref{fig:mobile_camera}所示，云端抓拍页面包含“刷新物模型”和“手动抓拍”两个按钮。页面会显示抓拍时间、触发原因和图片预览区域；在方案一验收阶段，设备只把BMP文件保存到SD卡，未开启图片上传，因此页面显示“暂无图片，请先抓拍或点刷新”。该提示与当前系统边界一致，不把App在线图片显示作为方案一通过条件。

\begin{figure}[H]
\centering
\thesisfig{抓拍.jpg}{0.42}
\caption{云端抓拍页面}
\label{fig:mobile_camera}
\end{figure}
\section{云端通信与核心功能}
通信模块负责与OneNET服务器建立数据连接，完成设备状态查询和控制命令下发。系统通过轮询方式读取设备当前属性，再把结果映射到移动端页面。设备侧根据物模型字段完成状态上报和控制响应。

本系统通信链路以ESP-01S（ESP8266）为WiFi接入模块，主控通过USART3与其通信，连线关系为PB10$\rightarrow$ESP\_RX、PB11$\leftarrow$ESP\_TX，模块使能脚CH\_PD/EN由PC9控制。默认工作波特率为115200\,bps，联网方式为STA模式，SSID与口令由统一宏定义管理，便于在不同网络环境下快速切换。该配置与第3章硬件连接保持一致，满足设备侧低成本联网和稳定会话维护需求\cite{esp8266_at,mqtt_spec,mqtt_v5}。

在初始化流程上，通信模块按“AT连通性检测$\rightarrow$入网$\rightarrow$MQTT会话建立$\rightarrow$主题订阅$\rightarrow$状态同步”顺序运行。为降低联调阶段故障定位成本，系统在启动阶段保留关键AT状态输出（固件版本、入网状态、连接结果），用于排查串口接线、波特率不匹配和入网异常等问题；在运行阶段通过重连策略和命令幂等机制维持弱网条件下的链路可用性。

在安全机制上，移动端到云平台接口采用HTTPS/TLS，设备端采用Token鉴权接入OneNET MQTT TCP链路，完成当前版本状态上报与远程控制闭环；设备侧MQTT TLS作为后续增强能力预留。总体来看，该通信模块形成了“移动端查询与控制—云平台转发—设备执行与回传”的完整路径，并与本课题“本地优先、远程协同”的系统目标保持一致。

移动端核心功能不再拆成多个短小三级标题，而按实际操作路径归纳为四项：  
\kuan{1}远程控制：用户可通过移动端对LED灯、布防状态和门锁状态进行控制。控制指令经云平台下发至设备端后，由设备完成相应操作，再把执行结果回传到平台，形成“发起--执行--确认”的闭环。  
\kuan{2}状态监控：页面实时展示温度、湿度、燃气浓度、告警状态、布防状态和门锁状态等关键数据。系统采用周期刷新机制更新页面，保证显示状态尽量和设备当前状态一致。  
\kuan{3}抓拍指令入口：当用户点击“立即抓拍”按钮后，移动端向云平台发送抓拍指令。设备端收到指令后启动本地抓拍流程，并优先把BMP文件写入SD卡。当前默认方案不要求立即回写图片地址，后续开启HTTP图床或MQTT分包后，再将图片地址与抓拍时间接入物模型字段。  
\kuan{4}图片查看预留：摄像头页面用于后续展示最新抓拍图片，并提供图片预览、复制链接等操作。当前阶段页面保留HTTP与MQTT两类预设访问位置，便于方案二接入后继续扩展。

\section{功能展示与实物场景}
根据本课题当前实现结果，系统已完成门禁控制、安防联动、环境异常处置、云端交互与稳定运行验证等核心能力。本节以文字归纳功能要点；对远程闭环、告警时序、断网重连、长稳运行等难以用单张静态照片完整呈现的过程，答辩与验收环节以录屏/实拍视频补充展示，论文中不再堆砌占位截图。

\kuan{1}本地门禁控制：系统已实现密码开锁与指纹开锁双通道，本地验证通过后执行开锁动作并记录状态，支持管理员功能窗口与本地配置管理。  
\kuan{2}远程控制闭环：系统通过OneNET与MQTT实现远程开锁和状态回传，设备侧可完成命令解析、执行反馈与属性上报。  
\kuan{3}入侵告警联动：在布防状态下，PIR触发可拉起告警流程，联动蜂鸣器、状态上报与事件记录。  
\kuan{4}燃气应急处置：MQ-2链路可完成阈值判定、关阀联动与持续告警，支持滤波与温漂补偿策略。  
\kuan{5}异常与恢复机制：系统具备断网降级、重连恢复、任务监测与看门狗保护能力，可在异常工况下维持关键功能。  
\kuan{6}稳定性验证：系统已完成分层测试与持续运行验证，关键指标满足既定设计目标。  

图\ref{fig:real_scene_prototype}给出系统实物搭建与室内场景模拟示意。图中重点标注PIR、DHT11、MQ-2、摄像头、OLED、键盘和指纹模块等功能部件，便于直接对应第3章整机电路图和第6章功能测试过程。上层微型家居布置只用于模拟室内空间，真正参与测试和评价的是箱体上的传感器、执行器、显示与交互模块。

\begin{figure}[H]
\centering
\thesisfigfit{images/fig5_4_real_scene_prototype_annotated.png}{0.90}{0.58}
\caption{实物搭建与真实场景模拟}
\label{fig:real_scene_prototype}
\end{figure}

\section{本章小结}
本章围绕\texttt{uni-app}移动端交互设计展开，完成了首页控制中心和摄像头抓拍页面设计，并通过OneNET云平台实现设备状态查询和控制命令下发。移动端可以显示关键环境数据，也支持远程控制。对于远程闭环、告警时序和断网重连这类不适合用单张图片展示的过程，本文以答辩视频演示作为补充，基本满足智能家居安防系统的人机交互需求。

\chapter{系统测试与结果分析}
\section{测试环境与方案}
测试方案按照工程验收习惯分层展开：先验证单模块能否稳定工作，再检查联动链路是否闭环，然后统计性能指标，最后进行长时运行和异常工况冲击\cite{freertos_book,stm32_ref}。这种安排可以逐层收敛问题，避免到系统级测试阶段才发现基础模块故障，从而增加定位成本。
测试环境与场景布置如图\ref{fig:test_setup}所示。

\begin{figure}[H]
\centering
\thesisfig{images/fig6_1_test_setup.png}{0.82}
\caption{测试环境与场景布置}
\label{fig:test_setup}
\end{figure}

OLED用于现场观察系统运行状态，其四行内容分别对应网络与心跳状态、门禁输入状态、温湿度/燃气采样值以及布防和门锁状态。图\ref{fig:oled_status_annotated}对实测界面中的文字含义进行了标注，便于在后续功能测试中把屏幕显示与系统状态对应起来。

\begin{figure}[H]
\centering
\thesisfig{images/fig6_0_oled_status_annotated.png}{0.48}
\caption{OLED状态显示字段含义标注}
\label{fig:oled_status_annotated}
\end{figure}

功能测试过程要求尽量把“操作动作”和“显示结果”放在同一证据链中记录。以DHT11湿度响应为例，测试时先记录常态读数，再将湿纸巾靠近传感器探头区域约10\,s，观察OLED上湿度值变化；停止干预后继续观察读数回落。DHT11湿度测试过程如图\ref{fig:dht11_humidity_wet_test}所示，键盘输入密码并完成本地开锁的测试过程如图\ref{fig:keyboard_unlock_test}所示，指纹录入测试过程如图\ref{fig:fingerprint_enroll_test}所示，MQ-2燃气超限报警测试过程如图\ref{fig:mq2_alarm_test}所示，远程开关灯测试过程如图\ref{fig:remote_light_test}所示，安防触发抓拍并保存到SD卡的过程如图\ref{fig:security_capture_sd_test}所示。

\begin{figure}[H]
\centering
\thesisfig{images/fig6_dht11_humidity_wet_test.png}{0.92}
\caption{DHT11湿度响应测试}
\label{fig:dht11_humidity_wet_test}
\end{figure}

\begin{figure}[H]
\centering
\thesisfig{images/fig6_keyboard_unlock_test.png}{0.92}
\caption{键盘输入密码后开锁测试}
\label{fig:keyboard_unlock_test}
\end{figure}

\begin{figure}[H]
\centering
\thesisfig{images/fig6_fingerprint_enroll_test.png}{0.94}
\caption{指纹录入测试过程}
\label{fig:fingerprint_enroll_test}
\end{figure}

\begin{figure}[H]
\centering
\thesisfig{images/fig6_mq2_alarm_test.png}{0.94}
\caption{MQ-2燃气超限报警测试}
\label{fig:mq2_alarm_test}
\end{figure}

\begin{figure}[H]
\centering
\thesisfig{images/fig6_remote_light_test.png}{0.92}
\caption{远程开关灯测试}
\label{fig:remote_light_test}
\end{figure}

\begin{figure}[H]
\centering
\thesisfig{images/fig6_security_capture_sd_test.png}{0.92}
\caption{安防触发抓拍并保存到SD卡测试}
\label{fig:security_capture_sd_test}
\end{figure}

为使测试结果可复现，测试统一在室内温度$25\pm2^\circ C$、相对湿度$50\%\pm10\%$条件下进行。网络分别设置为稳定Wi-Fi场景（平均RSSI约$-52$\,dBm）和波动场景（RSSI在$-72\sim-60$\,dBm变化）。除特别说明外，各项测试重复执行不少于30次，统计指标采用均值、标准偏差和95\%分位（P95）。测试环境关键参数如表\ref{tab:test_env_params}。

\begin{table}[htbp]
\centering
\caption{测试环境配置参数}
\label{tab:test_env_params}
\renewcommand{\arraystretch}{1.15}
\begin{tabular}{>{\raggedright\arraybackslash}p{3.2cm} >{\raggedright\arraybackslash}p{4.2cm} >{\raggedright\arraybackslash}p{4.8cm}}
\toprule
类别 & 配置参数 & 说明 \\
\midrule
硬件环境 & STM32F103RCT6主控板、AS608、HC-SR501、DHT11、MQ-2、OV7670+FIFO、继电器执行器 & 与第3章硬件设计一致，固件版本固定 \\
软件环境 & Keil MDK 5.x，FreeRTOS V9.x，OneNET MQTT 3.1.1，\texttt{uni-app}移动端 & 禁止测试中动态改库版本 \\
测试工具 & 可编程直流电源、数字万用表、示波器、逻辑分析仪、串口助手、网络抓包工具 & 用于时序、电压、电流、链路和日志采集 \\
测试条件 & 温度$25\pm2^\circ C$，湿度$50\%\pm10\%$，220\,V市电，稳定与波动网络双场景 & 异常工况测试在第6.6节单独执行 \\
\bottomrule
\end{tabular}
\end{table}
\section{单元功能测试}
单元测试阶段主要确认一件事：每个模块在独立条件下能否稳定完成自身功能。测试对象包括指纹、PIR、DHT11温湿度采样、MQ-2、执行器和串口通信。先让单元层通过，再进入联动层，可以减少后续联调中的连锁误判。对应标准化用例如表\ref{tab:unit_cases}，统计结果如表\ref{tab:unit_results}。

{\zihao{5}\renewcommand{\arraystretch}{1.2}\begingroup\setlength{\tabcolsep}{4.2pt}%
\begin{longtable}{>{\raggedright\arraybackslash}p{1.2cm} >{\raggedright\arraybackslash}p{1.9cm} >{\raggedright\arraybackslash}p{2.1cm} >{\raggedright\arraybackslash}p{2.2cm} >{\raggedright\arraybackslash}p{2.2cm} >{\raggedright\arraybackslash}p{2.2cm} >{\raggedright\arraybackslash}p{1.8cm}}
\caption{单元功能标准化测试用例}\\
\label{tab:unit_cases}\\
\toprule
用例ID & 测试目的 & 测试环境 & 测试步骤 & 预期结果 & 实际结果 & 判定标准 \\
\midrule
\endfirsthead
\multicolumn{7}{c}{表~\thetable\quad 单元功能标准化测试用例（续）}\\
\toprule
用例ID & 测试目的 & 测试环境 & 测试步骤 & 预期结果 & 实际结果 & 判定标准 \\
\midrule
\endhead
\bottomrule
\endfoot
UT-01 & 验证AS608识别功能 & 室温$25\pm2^\circ C$，串口57600 & 连续录入并识别50次合法指纹 & 成功率$\geq95\%$ & 49/50，98.0\%，通过 & 成功率$\geq95\%$ \\
UT-02 & 验证PIR触发有效性 & 布防开启，触发距离约3\,m & 布防状态下人员通过触发50次 & 触发率$\geq95\%$ & 48/50，96.0\%，通过 & 触发率$\geq95\%$ \\
UT-03 & 验证DHT11采样稳定性 & 恒温箱$25^\circ C$，采样周期2\,s & 50次周期读取并比对范围 & 无超量程与超时错误 & 50/50，100\%，通过 & 连续50次无异常 \\
UT-04 & 验证MQ-2阈值检测 & 预热30\,min，标准烟源注入 & 50次模拟超阈值场景采样 & 告警触发率$\geq90\%$ & 47/50，94.0\%，通过 & 触发率$\geq90\%$ \\
UT-05 & 验证执行器驱动链路 & 门锁12\,V供电，蜂鸣器在线 & 门锁/蜂鸣器各触发50次 & 动作正确且无卡滞 & 50/50，100\%，通过 & 无失动作/卡滞 \\
UT-06 & 验证串口JSON输出完整性 & STM32-ESP8266链路，115200\,bps & 连续上报100帧并核验字段 & 丢帧率$\leq1\%$ & 100/100，100\%，通过 & 字段完整且无丢帧 \\
\end{longtable}\endgroup
}

\begin{table}[htbp]
\centering
\caption{单元功能测试结果}
\label{tab:unit_results}
\renewcommand{\arraystretch}{1.15}
\begin{tabular}{>{\raggedright\arraybackslash}p{3.2cm} >{\raggedright\arraybackslash}p{2.2cm} >{\raggedright\arraybackslash}p{2.6cm} >{\raggedright\arraybackslash}p{2.2cm} >{\raggedright\arraybackslash}p{3.0cm}}
\toprule
测试模块 & 测试次数 & 成功次数 & 成功率(\%) & 备注 \\
\midrule
AS608指纹识别 & 50 & 49 & 98.0 & 串口57600，识别稳定 \\
HC-SR501人体检测 & 50 & 48 & 96.0 & 布防模式下触发有效 \\
DHT11采样 & 50 & 50 & 100.0 & 单总线读取正常 \\
MQ-2气体检测 & 50 & 47 & 94.0 & 受预热与阈值影响 \\
门锁/蜂鸣器控制 & 50 & 50 & 100.0 & GPIO驱动稳定 \\
串口JSON输出 & 100 & 100 & 100.0 & 约2s周期输出 \\
\bottomrule
\end{tabular}
\end{table}
\section{系统联动测试}
联动测试重点验证多模块协同能力和异常事件闭环处理效果，确保系统在典型家庭场景下能够自动执行安全策略。联动测试用例和结果分别如表\ref{tab:linkage_cases}与表\ref{tab:linkage_results}。

{\zihao{5}\renewcommand{\arraystretch}{1.2}\begingroup\setlength{\tabcolsep}{4.2pt}%
\begin{longtable}{>{\raggedright\arraybackslash}p{1.2cm} >{\raggedright\arraybackslash}p{1.8cm} >{\raggedright\arraybackslash}p{2.2cm} >{\raggedright\arraybackslash}p{2.3cm} >{\raggedright\arraybackslash}p{2.1cm} >{\raggedright\arraybackslash}p{2.1cm} >{\raggedright\arraybackslash}p{1.8cm}}
\caption{联动测试标准化用例}\\
\label{tab:linkage_cases}\\
\toprule
用例ID & 测试目的 & 测试环境 & 测试步骤 & 预期结果 & 实际结果 & 判定标准 \\
\midrule
\endfirsthead
\multicolumn{7}{c}{\tablename\ \thetable\ （续）}\\
\toprule
用例ID & 测试目的 & 测试环境 & 测试步骤 & 预期结果 & 实际结果 & 判定标准 \\
\midrule
\endhead
\bottomrule
\endfoot
LT-01 & 门禁联动闭环 & 本地模式，日志记录开启 & 合法/非法指纹交替触发30轮 & 开锁正确，非法超阈值告警 & 平均1.20\,s，29/30通过 & 成功率$\geq95\%$ \\
LT-02 & 入侵告警联动 & 布防开启，摄像头在线 & PIR触发并检查本地抓拍落盘与告警上报 & 蜂鸣、抓拍落盘、状态上报均触发 & 平均1.35\,s，28/30通过 & 成功率$\geq90\%$ \\
LT-03 & 燃气应急联动 & MQ-2预热完成，阀门在线 & 模拟气体超限并观察关阀告警 & 关阀+告警+持续上报 & 平均1.10\,s，29/30通过 & 关阀动作100\% \\
LT-04 & 远程控制闭环 & OneNET在线，弱网波动场景 & 移动端下发开锁/布防命令30轮 & 执行成功并回执状态 & 平均1.90\,s，28/30通过 & 成功率$\geq90\%$ \\
LT-05 & 复合事件协同 & PIR+MQ-2+远程命令并发 & 并发触发并记录调度顺序 & 安全任务优先，闭环不中断 & 30/30通过 & 关键路径无死锁 \\
\end{longtable}\endgroup
}
\kuan{1}门禁联动：合法用户指纹匹配成功后，系统可在短时间内完成开锁和日志记录；非法尝试达到阈值时，系统触发本地告警并完成移动端状态更新。测试结果说明，门禁联动流程完整，执行过程也比较稳定。  
\kuan{2}入侵告警：布防状态下触发人体红外后，系统可完成告警、图像抓拍和消息上报。通过设置延时确认与状态复核机制，可有效降低短时误触发造成的误报率，也说明移动端和云平台之间的联动链路能够正常工作。  
\kuan{3}燃气应急：在模拟气体超限场景中，系统会自动执行切断燃气和蜂鸣器告警动作，并持续上报状态。测试结果验证了应急流程的可靠性和及时性，同时说明移动端在异常状态展示方面能够保持同步更新。  
\kuan{4}远程控制：用户通过移动端下发控制命令后，系统可正确执行并返回状态反馈。多次重复测试中，命令执行成功率保持在较高水平，满足远程监测和控制需求。

联动测试结果有两个比较明显的特征。第一，本地三条关键链路，也就是门禁、入侵和燃气应急，平均响应时间都在2\,s以内，说明本地闭环路径比较稳定。第二，远程控制时延更高，离散性也更大，主要受无线波动和重连窗口影响。回看异常样本后，没有发现状态机失稳或本地逻辑崩溃，问题主要集中在通信侧。因此，后续优化重点应放在连接保持、重传窗口和命令去重上，而不是反复改动核心控制逻辑。

\begin{table}[htbp]
\centering
\caption{系统联动测试结果}
\label{tab:linkage_results}
\renewcommand{\arraystretch}{1.15}
\begin{tabular}{>{\raggedright\arraybackslash}p{3.2cm} >{\raggedright\arraybackslash}p{2.0cm} >{\raggedright\arraybackslash}p{2.6cm} >{\raggedright\arraybackslash}p{2.4cm} >{\raggedright\arraybackslash}p{2.6cm}}
\toprule
测试场景 & 测试次数 & 正常完成次数 & 平均响应时间(s) & 异常说明 \\
\midrule
门禁联动 & 30 & 29 & 1.20 & 指纹成功后开锁 \\
入侵告警 & 30 & 28 & 1.35 & PIR触发后蜂鸣与抓拍 \\
燃气应急 & 30 & 29 & 1.10 & 超阈值后关阀告警 \\
远程控制 & 30 & 28 & 1.90 & 受网络波动影响 \\
\bottomrule
\end{tabular}
\end{table}
\section{性能指标测试}
性能测试从准确性、实时性、功耗与图像链路四个维度评估系统能力。数据处理流程为：原始日志采集$\rightarrow$异常值剔除（3$\sigma$规则）$\rightarrow$均值与标准差计算$\rightarrow$P95统计$\rightarrow$与目标值对比分析。设计目标与实测结果逐项对齐如表\ref{tab:perf_target_compare}，离散性统计如表\ref{tab:perf_stat}；图\ref{fig:accuracy_bar}和图\ref{fig:latency_line}用于可视化展示关键指标。
\kuan{1}准确率与误报率：本文通过多轮样本测试统计门禁识别准确率与入侵检测误报率。测试结果说明，系统在典型室内环境中具有较高识别准确性，误报率控制在可接受范围内，移动端显示的数据与设备端状态基本一致。  
\kuan{2}响应延迟：对门禁开锁、告警触发与远程控制等关键路径进行时延测量。结果显示本地事件响应快速，远程控制时延主要受网络环境影响，但整体满足家庭应用需求。  
\kuan{3}功耗：系统功耗测试分为空闲、监测与告警三种状态。结果显示空闲与常规监测状态功耗较低，告警状态因执行器与通信负载增加而上升，整体仍处于可控范围。  
\kuan{4}本地抓拍链路：图像链路测试关注触发采集成功率、本地写入耗时与BMP文件可用性。测试表明在当前硬件条件下，系统可完成OV7670+FIFO到SD卡BMP的本地抓拍流程。图片上传与移动端展示作为后续增强功能保留，不作为方案一的通过条件。

结合表\ref{tab:perf_target_compare}和表\ref{tab:perf_stat}可以得到两点实际结论。第一，本地闭环能力已经达到设计目标，关键功能在准确率、时延和功耗上都满足约束。第二，系统上限主要受网络质量和部分外设物理特性影响，比较典型的是MQ-2预热漂移、无线抖动和抓拍链路带宽占用。针对这些边界，后续应优先做参数层优化，例如阈值+趋势联合判定、分级重连、本地抓拍稳定优先，在不增加硬件成本的前提下提升稳态表现。
关键功能准确率统计如图\ref{fig:accuracy_bar}所示。

\begin{figure}[H]
\centering
\thesisfig{images/fig6_2_accuracy_bar.png}{0.78}
\caption{关键功能准确率统计}
\label{fig:accuracy_bar}
\end{figure}

关键路径响应时延对比如图\ref{fig:latency_line}所示。
\begin{figure}[H]
\centering
\thesisfig{images/fig6_3_latency_line.png}{0.78}
\caption{关键路径响应时延对比}
\label{fig:latency_line}
\end{figure}

\begin{table}[htbp]
\centering
\caption{关键性能指标测试数据}
\label{tab:perf_target_compare}
\renewcommand{\arraystretch}{1.15}
\begin{tabular}{>{\raggedright\arraybackslash}p{3.5cm} >{\raggedright\arraybackslash}p{2.8cm} >{\raggedright\arraybackslash}p{2.5cm} >{\raggedright\arraybackslash}p{2.6cm}}
\toprule
指标项 & 测试结果 & 目标值 & 是否达标 \\
\midrule
门禁识别准确率(\%) & 98.0 & $\geq$95.0 & 是 \\
入侵检测误报率(\%) & 3.2 & $\leq$5.0 & 是 \\
本地告警响应时延(ms) & 620 & $\leq$1000 & 是 \\
远程控制响应时延(ms) & 1850 & $\leq$2500 & 是 \\
空闲功耗(mW) & 420 & $\leq$600 & 是 \\
告警状态功耗(mW) & 980 & $\leq$1200 & 是 \\
本地抓拍写入成功率(\%) & 96.5 & $\geq$90.0 & 是 \\
\bottomrule
\end{tabular}
\end{table}
\begin{table}[htbp]
\centering
\caption{关键性能统计结果（含离散性）}
\label{tab:perf_stat}
\renewcommand{\arraystretch}{1.15}
\begin{tabular}{>{\raggedright\arraybackslash}p{3.3cm} >{\centering\arraybackslash}p{2.0cm} >{\centering\arraybackslash}p{2.2cm} >{\centering\arraybackslash}p{2.0cm} >{\centering\arraybackslash}p{2.3cm}}
\toprule
指标项 & 样本数 & 均值 & 标准差 & P95 \\
\midrule
本地告警响应时延(ms) & 120 & 620 & 84 & 790 \\
远程控制响应时延(ms) & 120 & 1850 & 260 & 2280 \\
本地抓拍写入耗时(ms) & 90 & 2110 & 330 & 2680 \\
空闲功耗(mW) & 60 & 420 & 18 & 455 \\
告警功耗(mW) & 60 & 980 & 42 & 1050 \\
\bottomrule
\end{tabular}
\end{table}
\section{稳定性测试}
稳定性测试采用连续运行和多事件压力触发方式，观察系统是否出现死机、丢包或任务阻塞。结果说明，系统在长时间运行中保持较好稳定性，串口状态输出、MQTT链路与任务调度机制运行正常。
72h稳定性测试流程如图\ref{fig:stability_flow}所示。

\begin{figure}[H]
\centering
\thesisfigfit{images/fig6_4_stability_flow.png}{0.92}{0.76}
\caption{72h稳定性测试流程图}
\label{fig:stability_flow}
\end{figure}

为进一步评估系统在复杂负载下的持续运行能力，测试阶段还对“高频状态上报+多事件连续触发+远程指令插入”复合场景进行观测。结果显示，系统在任务优先级配置合理的前提下能够维持核心安防任务实时性，低优先级非关键任务可能出现轻微延后但不影响安全闭环。该现象与预期一致，说明当前调度策略具备“关键路径优先保障”特性，符合家庭安防系统设计目标。

\begin{table}[htbp]
\centering
\caption{稳定性测试记录}
\label{tab:stability_record}
\renewcommand{\arraystretch}{1.15}
\begin{tabular}{>{\raggedright\arraybackslash}p{3.0cm} >{\raggedright\arraybackslash}p{2.5cm} >{\raggedright\arraybackslash}p{2.5cm} >{\raggedright\arraybackslash}p{2.5cm} >{\raggedright\arraybackslash}p{2.5cm}}
\toprule
运行时长(h) & 重启次数 & 通信异常次数 & 任务阻塞次数 & 结论 \\
\midrule
24 & 0 & 1 & 0 & 稳定 \\
48 & 0 & 2 & 0 & 稳定 \\
72 & 0 & 3 & 0 & 稳定 \\
\bottomrule
\end{tabular}
\end{table}
\section{异常工况与鲁棒性测试}
为验证系统工程鲁棒性，增加断网离线、传感器故障、电源波动和强电磁干扰四类异常工况测试，测试用例与结果如表\ref{tab:abnormal_cases}。

{\zihao{5}\renewcommand{\arraystretch}{1.2}
\begin{longtable}{>{\raggedright\arraybackslash}p{1.2cm} >{\raggedright\arraybackslash}p{1.8cm} >{\raggedright\arraybackslash}p{2.6cm} >{\raggedright\arraybackslash}p{2.1cm} >{\raggedright\arraybackslash}p{2.3cm} >{\raggedright\arraybackslash}p{2.0cm}}
\caption{异常工况测试用例与结果}\\
\label{tab:abnormal_cases}\\
\toprule
用例ID & 异常类型 & 测试方法 & 预期结果 & 实测结果 & 系统容错表现 \\
\midrule
\endfirsthead
\multicolumn{6}{c}{\tablename\ \thetable\ （续）}\\
\toprule
用例ID & 异常类型 & 测试方法 & 预期结果 & 实测结果 & 系统容错表现 \\
\midrule
\endhead
\bottomrule
\endfoot
AT-01 & 断网离线 & 运行中断开Wi-Fi 10\,min后恢复，重复20轮 & 本地告警不中断，恢复后补传日志 & 20/20无漏告警，补传成功率100\% & 离线自治与队列补传有效 \\
AT-02 & 传感器故障 & 拔除DHT11/MQ-2信号线并保持5\,min，重复20轮 & 故障上报并降级，其他链路可用 & 20/20触发故障标志，主循环无阻塞 & 任务隔离与故障降级有效 \\
AT-03 & 电源波动 & 施加$\pm10\%$波动及100\,ms瞬降，重复30轮 & 不死机，关键任务不中断，可自动恢复 & 30/30无异常复位，记录2次欠压事件 & 电源监测与看门狗协同有效 \\
AT-04 & 强电磁干扰 & 继电器邻近快速开关并叠加对讲机近距离发射，重复30轮 & 无持续误告警，通信可恢复 & 30轮中瞬态误触发1次，3\,s内恢复 & 去抖与多次确认策略有效 \\
\end{longtable}
}
\section{测试结论与优化}
本系统对MQ-2气体传感器完成了标准化实验室标定。结合第3章表\ref{tab:mq2_calib_raw}可得：MQ-2通过多浓度点标定建立了ADC读数与气体浓度等级的映射关系，可稳定实现分级检测与告警\cite{mq2_calib_source}。同时，DHT11链路在单元测试与异常工况测试中均保持稳定，能够支撑环境状态展示和联动辅助判定。

综合联动测试和长期稳定性测试结果可知，环境采集链路整体满足系统运行需求。后续优化重点仍聚焦于误触发抑制、弱网重传效率、低功耗策略与图片上传增强链路，以进一步提升整机工程可用性与用户体验。
\section{误差来源与可信度分析}
为增强测试结论的可信度，本文对误差来源进行了归因分析。系统误差主要来自三类因素：传感器物理特性、采样时序和网络环境。具体来看，MQ-2与PIR对温湿度、预热状态和安装角度比较敏感，可能带来阈值漂移；任务切换与中断触发会引入毫秒级采样抖动，对边界样本判定有一定影响；远程控制链路还会受无线信道质量和平台转发负载影响，导致时延分布出现长尾。这些误差很难完全消除，但可以通过合理实验设计把它们对结论有效性的影响控制在较低水平。
主要误差来源及其作用关系如图\ref{fig:error_fishbone}所示。
\begin{figure}[H]
\centering
\thesisfigfit{images/fig6_5_error_fishbone.png}{0.86}{0.56}
\caption{测试误差来源分析图}
\label{fig:error_fishbone}
\end{figure}

\kuan{1}重复试验与样本代表性：为避免“偶然成功”被当作稳定能力，本文对多数测试项设置了不少于30次重复，核心模块扩大到50次或100次。该规模足以支撑本科毕设层面的结论判断，但样本仍主要来自常温室内和中等网络波动环境。因而本文结论的适用范围应限定在已测试场景内，对极端温度、高密无线和强干扰边界的外推仍需后续补测。  
\kuan{2}统计指标选择与局限：当前统计口径以成功率、平均时延和稳定运行记录为主，便于与设计指标直接对照；不足是均值对长尾敏感性有限。本文已补充标准差和P95用于描述离散性，但仍可继续加入中位数和分场景统计，以更完整反映时延分布。对于误报率这类概率指标，后续建议扩大跨时段样本并按场景分层，减少单一环境偏置。  
\kuan{3}结果可信度提升策略：结合本项目实施过程，提升可信度最有效的做法有三项：固定测试基线、完整保留过程数据、对异常样本做回放复盘。固定基线用于保证前后版本可比；过程数据（日志、时间戳、状态快照）用于定位误差来源；异常复盘用于区分系统性缺陷和偶发扰动。按这套流程执行后，测试结论可追溯性更好，也便于后续版本按证据迭代。

本系统误差修正措施严格对应前文的校准、滤波与信号调理实现，并以可量化指标作为验收口径：  
\kuan{1}MQ-2温漂与噪声误差修正（对应式~\eqref{eq:mq2_temp_comp}、表\ref{tab:mq2_filter_compare}）：对基线ADC引入温漂补偿$A_{\mathrm{comp}}=A_{\mathrm{raw}}-\alpha_T(T-25)$（$\alpha_T=3.2$\,ADC/$^\circ$C），使15/25/35$^\circ$C基线由172/205/239补偿收敛至约204/205/207；同时采用“中值滤波+滑动平均”级联（式~\eqref{eq:median_filter}、式~\eqref{eq:mean_filter}），抑制继电器动作尖峰，实测标准差由约31.4\,ADC降至约9.7\,ADC（下降约69.1\%），显著降低误触发概率。  
\kuan{2}阈值与判据设定（对应第3章表\ref{tab:mq2_calib_raw}与式~\eqref{eq:mq2_temp_comp}）：MQ-2阈值采用“基线均值+3$\sigma$裕量+安全余量”原则，取$A_{\mathrm{base}}=205$、$\sigma=18$、$\Delta_{\mathrm{safe}}=40$得到$T_{gas}\approx299$（工程取300）；并与变化率阈值$\Delta T$联合形成“双判据”（$\bar{A}_k>T_{gas}$且$\Delta A_k>\Delta T$），避免短时干扰仅凭单阈值触发。  
\kuan{3}硬件链路抑噪（对应式~\eqref{eq:decouple}、式~\eqref{eq:decouple_parallel}、式~\eqref{eq:rc_mag}）：主控去耦按并联总容值$C_{\mathrm{sum}}=n\times0.1\,\mu$F选取$n=6$满足0.6\,$\mu$F瞬态需求，并叠加10\,$\mu$F储能；模拟采样支路采用分压+RC低通，其幅频响应见式~\eqref{eq:rc_mag}，在20\,kHz/50\,kHz分别衰减约$-22.0$/$-29.9$\,dB，实测ADC输入纹波由约120\,mV降至约34\,mV（降噪约71.7\%），对应第6章异常工况与稳定性测试中继电器动作引起的采样尖峰显著减少。  
\kuan{4}链路时延与长尾抑制（对应第6章表\ref{tab:perf_stat}与置信区间分析）：对远程控制采用幂等去重（\texttt{cmdId+ts}）与分级重传，降低弱网下重复执行风险；时延统计采用均值/标准差/P95并与目标对齐，保证结论可复现可追溯。  
以上修正措施均可在现有硬件条件下通过固件配置或算法参数更新完成，形成与论文第3章（调理与校准）—第4章（滤波与判定）—第6章（测试验证）一致的工程闭环。

\section{本章小结}
本章搭建了可复现的测试环境，采用标准化测试用例完成单元、联动、性能、稳定性和异常工况验证，并用统计指标、曲线和柱状图对结果进行量化分析。测试结果说明，系统在断网、传感器异常、电源波动和并发压力场景下仍能保持核心功能可用，关键性能指标均达到设计目标。通过误差归因和修正措施，本文形成了面向测控工程的验证闭环，也为后续版本优化提供了依据。

\chapter{总结与展望}
\section{总结}
本文围绕家庭安防与智能家居联动需求，设计并实现了一套以STM32F103RCT6为核心的单芯片安防联动系统。系统完成了密码/指纹门禁、PIR入侵检测、MQ-2燃气异常处理、DHT11温湿度采集、继电器与蜂鸣器联动、OV7670+FIFO本地抓拍落盘、OneNET状态上报和移动端远程控制等功能。整体实现没有把关键安全动作完全依赖云端，而是让开锁、关阀、蜂鸣和本地抓拍优先在设备端完成，再通过网络进行状态同步和远程协同。

测试结果说明，系统门禁识别准确率达到98.0\%，入侵误报率为3.2\%，本地告警平均时延为620\,ms，远程控制平均时延为1850\,ms，本地抓拍写入成功率为96.5\%，72\,h连续运行未出现死机和任务阻塞。图像链路按当前工程边界采用分阶段方案：方案一完成本地SD抓拍和状态闭环，HTTP图床或MQTT分包上传作为后续增强。上述结果说明，本系统在本科毕业设计条件下能够完成多模块协同，并具备可实现、可验证和可复现的工程特点。
\section{各章节工作对应总结}
第1章从家庭安防实际需求出发，分析了智能家居联网、居住场景安全事件和弱网条件下本地闭环的重要性，并结合国内外研究现状确定了“本地自治+远程协同”的技术主线。第2章根据样机实际功能完成需求分解，明确门禁、入侵、燃气应急、远程协同和本地抓拍等闭环任务，并给出总体架构、单芯片控制策略、系统流程和关键技术选型。

第3章完成硬件设计与整机接线说明，内容包括STM32F103RCT6最小系统、AS608、HC-SR501、DHT11、MQ-2、OV7670+FIFO、继电器、蜂鸣器、ESP8266、SD卡和电源模块等。该章还给出整机原理图、接口分配、电源预算和硬件联调案例，为后续软件任务和测试验证提供硬件基础。

第4章完成软件架构和核心逻辑设计。系统采用裸机联调与FreeRTOS运行双模式，围绕门禁验证、安防告警、燃气应急、场景联动、MQTT通信、异常恢复和看门狗保护建立软件闭环。第5章完成移动端交互设计，首页用于显示温湿度、燃气、门锁、布防、报警和阈值配置，云端抓拍页用于发送抓拍命令和预留图片展示位置。第6章通过单元测试、系统联动测试、性能测试、稳定性测试和异常工况测试，对前文设计目标进行了验证。
\section{创新点与优势}
结合本课题的实际开发和测试过程，本文的创新点与优势主要体现在以下几个方面：  
\kuan{1}单芯片集中控制的本地闭环架构。系统使用STM32F103RCT6统一管理门禁、传感器、执行器、抓拍和联网模块，减少了多控制器同步和板间通信带来的状态冲突。测试中门禁识别准确率达到98.0\%，本地告警平均时延为620\,ms，说明该架构可以满足家庭安防的实时性要求。  
\kuan{2}关键安全动作本地优先。门禁开锁、燃气阀切断、蜂鸣告警和本地抓拍不等待云端确认即可执行，OneNET和移动端主要承担状态显示与远程协同。这样在弱网或短时断网时，系统仍能保留基础安防能力。  
\kuan{3}多传感复核降低误报。入侵链路不是简单把PIR高电平直接判定为告警，而是加入持续时间、锁定窗口和抓拍复核等条件。第6章统计中入侵误报率为3.2\%，满足设计目标。  
\kuan{4}面向实际噪声的信号处理。MQ-2通道结合温漂补偿、中值滤波和滑动平均处理，降低继电器动作、电源纹波和瞬态干扰对ADC结果的影响。该处理链实现复杂度不高，适合在STM32平台上落地。  
\kuan{5}测试口径清晰，便于复现。本文将门禁准确率、入侵误报率、本地告警时延、远程控制时延、本地抓拍写入成功率和72\,h稳定性作为主要验收指标，并在第6章逐项测试。这样的验证方式避免只停留在功能演示层面，也便于后续版本继续对比。

总体来看，本文在保持前文结论不变的基础上，对毕设工作进行了面向工程落地的总结和复盘。全文完成了从需求分析、方案设计、硬件实现、软件协同、平台交互到测试验证的完整闭环，验证了单芯片家用安防联动方案在本科毕业设计尺度下的可行性和应用价值。后续工作可在现有架构优势的基础上，继续向更高智能化、更强鲁棒性和更好可维护性方向迭代。

\clearpage
\begin{center}
{\heiti\zihao{3}参考文献}
\end{center}
\phantomsection
\addcontentsline{toc}{chapter}{参考文献}
\printbibliography[heading=none]

\chapter*{致谢}
\addcontentsline{toc}{chapter}{致谢}
时光匆匆，四年广东工业大学测控技术与仪器专业本科求学之旅转瞬落幕。从初入校园的专业懵懂，到深耕嵌入式智能测控技术，随着《家用安防与智能家居联动系统设计》毕业设计顺利完成，我的本科学业也画上圆满句号。回首求学时光，专业知识的沉淀、毕业设计孤身攻坚调试、独自备战本校仪器仪表工程研究生复试的点滴成长，皆离不开师长教诲与家人扶持，在此我致以最诚挚的谢意。

衷心感谢我的毕业设计指导老师。本次毕设从课题选题、系统整体架构规划，到STM32F103RCT6单芯片核心方案优化，以及ESP8266、各类传感器、摄像头外设调试和程序代码BUG修复，全程均由我独立研发实操。面对硬件接线故障、代码逻辑冲突、系统联动不稳等诸多实操难题，老师始终悉心指导我梳理调试思路，传授严谨工程化研发思维。在论文撰写修改阶段，老师细致打磨文稿内容、规范学术格式，给出专业中肯的修改建议，助力我毕业设计顺利落地、论文完善成型。老师求真务实的科研态度与精益求精的治学精神，为我后续研究生科研深造树立了良好榜样，让我受益良多。

感谢广东工业大学的悉心培养，感谢测控专业全体授课老师。各位老师课堂倾囊相授测控原理、嵌入式开发、传感器检测等核心专业知识，为我筑牢理论根基；课余耐心答疑解惑，指引我将课本理论与工程实操相结合。依托学校优良学风和完善实训条件，我从零自学嵌入式开发，积累项目实战能力，也为研究生复试备考深造打下坚实专业基础。

感恩我的家人，你们是我求学路上最坚实的后盾。无论是独自熬夜攻坚毕设调试工作，还是全心投入研究生复试备考，你们始终默默包容我的忙碌，在生活上悉心照料、精神上全力支持，让我能够心无旁骛专注学业与追梦前行，毫无后顾之忧奔赴科研深造理想，养育之恩，铭记于心。

最后，感谢始终坚持、从未轻言放弃的自己。本科四年，自学钻研专业知识，独自打磨毕设项目、反复调试攻坚克难，孤身备战研究生复试。一路走来，虽有调试受挫、备考迷茫，但终在坚持中收获成长，不仅提升了专业实操能力，更磨砺了心性与毅力。

本科征程落幕，新的科研前路启程。未来我将怀揣感恩之心，秉持广工治学精神，深耕仪器仪表与智能嵌入式领域，脚踏实地潜心钻研，不负教诲、不负韶华、不负初心。

\appendix
\chapter{附录}
\section{系统原理图}
本课题的原理图、接线图和整机布局图已集中放置在第3章硬件设计部分，避免同一张电路图在正文和附录中重复出现。附录中保留归档说明，便于后续查阅和替换源文件。主要图纸包括主控最小系统、AS608指纹模块、HC-SR501人体红外、DHT11单总线、OV7670+FIFO图像采集链路、蜂鸣器驱动、ESP8266联网模块、MicroSD模块、电源管理模块、整机电路布局与整机原理图。上述图纸与第3章表\ref{tab:power_budget}和接口分配表共同构成硬件设计归档。

其中，图\ref{fig:whole_circuit_layout}用于说明STM32核心板与传感、执行、通信、存储和电源模块之间的总体接线关系；图\ref{fig:system_overall_schematic}给出整机原理图，用于核对各模块排针、信号名、电源和共地关系。

\section{核心代码}
本节给出系统实现中与论文功能闭环直接相关的核心代码截图。代码选取原则为：优先覆盖任务调度、门禁与安防联动、MQ-2燃气应急、云端命令解析、本地抓拍落盘、看门狗恢复和下行命令幂等处理等关键路径；不重复展示普通初始化、宏定义和第三方库代码。

\subsection*{代码片段1：任务启动与调度配置（节选自 app\_rtos.c）}
\begin{figure}[H]
\centering
\thesisfigfit{images/appendix_code_1_main_loop.png}{0.96}{0.52}
\caption{FreeRTOS任务启动与调度配置核心代码}
\end{figure}

\subsection*{代码片段2：门禁与安防联动处理（节选自 lock\_manager.c / linkage.c）}
\begin{figure}[H]
\centering
\thesisfigfit{images/appendix_code_2_security_alarm.png}{0.96}{0.46}
\caption{门禁状态机与安防联动核心代码}
\end{figure}

\subsection*{代码片段3：MQ-2燃气应急与抓拍流程（节选自 linkage.c / Capture.c）}
\begin{figure}[H]
\centering
\thesisfigfit{images/appendix_code_3_mq2_emergency.png}{0.96}{0.44}
\caption{MQ-2燃气应急与本地抓拍核心代码}
\end{figure}

\subsection*{代码片段4：云端命令解析与执行（节选自 esp8266\_tls.c）}
\begin{figure}[H]
\centering
\thesisfigfit{images/appendix_code_4_cloud_cmd.png}{0.96}{0.44}
\caption{OneNET下行命令解析与执行核心代码}
\end{figure}

\subsection*{代码片段5：本地抓拍与落盘流程（节选自 Capture.c）}
\begin{figure}[H]
\centering
\thesisfigfit{images/appendix_code_5_capture_pipeline.png}{0.96}{0.46}
\caption{本地抓拍与BMP落盘核心代码}
\end{figure}

\subsection*{代码片段6：多源看门狗判定与喂狗（节选自 wdg.c）}
\begin{figure}[H]
\centering
\thesisfigfit{images/appendix_code_6_wdg_recovery.png}{0.96}{0.44}
\caption{多源看门狗判定与恢复核心代码}
\end{figure}

\subsection*{代码片段7：下行命令幂等去重（节选自 esp8266\_tls.c）}
\begin{figure}[H]
\centering
\thesisfigfit{images/appendix_code_7_cmd_idempotent.png}{0.96}{0.42}
\caption{云端下行命令幂等去重核心代码}
\end{figure}

\section{实物图}
实物与测试图片已按使用场景放入正文对应章节：第5章保留系统实物搭建与真实场景模拟图，用于说明各模块在样机箱体中的空间布置；第6章集中放置功能测试过程图，用于证明“操作动作—显示结果—恢复/闭环结果”的证据链。当前已归档的实测图片包括OLED状态显示字段含义标注、DHT11湿度响应测试、键盘输入密码后开锁测试、指纹录入测试过程、MQ-2燃气超限报警测试、远程开关灯测试，以及安防触发抓拍并保存到SD卡测试。

\end{document}
